%
% 東京電機大学
% システムデザイン工学研究科
% 卒業論文アウトライン
%
\documentclass[twocolumn,10pt,a4paper]{jsarticle}

% ページ設定
\textheight 255mm
\textwidth 170mm
\topmargin -17mm
\oddsidemargin -5mm
\evensidemargin 0mm
\columnsep 8mm

% section,subsection,subsubsectionのフォントサイズの設定
\usepackage{titlesec}
\titleformat*{\section}{\large\bfseries}
\titleformat*{\subsection}{\normalsize\bfseries}
\titleformat*{\subsubsection}{\small\bfseries}

% 図の設定
\usepackage[dvipdfmx]{graphicx}
\renewcommand{\figurename}{Fig.} % 図->Fig.

% 表の設定
\renewcommand{\tablename}{Table} % 表->Table

% 概要の見出しの設定
\renewcommand{\abstractname}{\bf ABSTRACT} % 概要->ABSTRACT

% 数学記号用パッケージ
\usepackage{amsmath}
\usepackage{amssymb}
\usepackage{tabularx}
\usepackage{booktabs}
\usepackage{float}

% TikZ 図作成用パッケージ
\usepackage{tikz}
\usetikzlibrary{positioning,arrows.meta,backgrounds,fit,calc}

% URL とハイパーリンク用パッケージ
\usepackage{url}
\usepackage[dvipdfmx]{hyperref}
% URL の改行を柔軟に
\def\UrlBreaks{\do\-\do\.\do\@\do\\\do\/\do\!\do\_\do\|\do\;\do\>\do\]%
\do\)\do\,\do\?\do\&\do\'\do+\do\=\do\#}
\Urlmuskip=0mu plus 1mu

% maketitleの設定(jsarticle用)
\makeatletter
\def\etitle#1{\def\@etitle{#1}}
\def\id#1{\def\@id{#1}}
\def\teacher#1{\def\@teacher{#1}}

\def\@maketitle{
\begin{center}
\fontsize{14pt}{18pt}\selectfont%
{\bf \@title \par}% 日本語タイトル
\fontsize{12pt}{16pt}\selectfont%
{\bf \@etitle \par}% 英語タイトル
\end{center}
\begin{flushright}
\fontsize{11pt}{15pt}\selectfont%
\begin{tabular}[t]{ll}%
学籍番号 & 22AJ603 \\
学生氏名 & ロ メイトウ \\
研究指導教員 & 朱 金暁 先生 \\
\end{tabular}%
\par\vskip 1.5em
\end{flushright}
\ifvoid\@abstractbox\else\centerline{\box\@abstractbox}\vskip1.5em\fi
}
\makeatother

% タイトル情報
\title{2.4\,GHz帯無線リモコンシステムを対象とした\\リプレイ攻撃防御手法の検討と評価}
\etitle{Evaluation of Replay Attack Defense Mechanisms\\
for 2.4\,GHz Wireless Remote Control Systems}
% 著者・教員情報は \@maketitle 内に直接記述
\date{}

\begin{document}

\baselineskip 12pt
\maketitle

\begin{abstract}
低価格な無線リモコンや IoT デバイスでは，コストや電力制約から暗号・認証機能が省略されることが多く，記録した無線フレームの単純な再送によるリプレイ攻撃が現実的な脅威となる．
本研究は，市販 2.4\,GHz 帯無線リモコン玩具カーを対象としてリプレイ脆弱性を明らかにするとともに，軽量なアンチリプレイ機能の有効性をチャネル特性を考慮して定量的に評価し，低コスト無線制御システムにおける設計指針を与えることを目的とする．

まず，ソフトウェア無線機と公開解析ツールを用いてリモコンフレームのビット列構造を解析し，HackRF One によるリプレイ実験により，正規リモコンを用いずに車体の走行や灯光を制御できることを示す．
次に，パケット損失率 $p_{\mathrm{loss}}$ と順序入れ替え率 $p_{\mathrm{reorder}}$ をもつ確率的チャネルモデルの下で，防御なし，ローリングカウンタ方式，スライディングウィンドウ方式，簡易チャレンジレスポンス方式の 4 方式を共通のモンテカルロシミュレータ上で比較し，正規受理率（LAR）と攻撃成功率（ASR）および総合スコアにより評価する．

数値実験（$p_{\mathrm{loss}} = 0.15$，$p_{\mathrm{reorder}} = 0.15$，inline 攻撃）では，受理ウィンドウ幅 $W$ を 3～7 としたスライディングウィンドウ方式が，可用性（LAR）を 85\% 前後に維持しつつ，攻撃成功率（ASR）を 6.5～8.7\% に抑える．
特に $W = 3$～$5$ は安全性と可用性の折衷が良い．
以上より，損失と乱序が混在する 2.4\,GHz 帯無線環境において，単純な状態管理のみでリプレイ攻撃を緩和するための実用的なアンチリプレイ設計指針を提示する．
\end{abstract}

%%%%%%%%%%%%%%%%%%%%%%%%%%%%%%%%%%%%%%%%%%%%%%%%
\section{はじめに}

\subsection{研究背景}

近年，センサ，家電，玩具，産業機器などをインターネットに接続する IoT（Internet of Things）が急速に普及している \cite{Sinha2025IoTAnalytics}．IoT デバイスの増加に伴い，無線通信を用いた低コスト機器が多様な場面で利用されるようになった \cite{Carpenter2025IEEE802154}．これらの多くは 2.4\,GHz ISM 帯（2.400--2.4835\,GHz）を利用し，Wi-Fi（IEEE 802.11），Bluetooth Low Energy，IEEE 802.15.4 系（ZigBee/Thread 等）を含む方式が広く用いられている \cite{BLE2025Sensors}．

一方で，低価格 IoT／無線リモコン機器では，コスト・消費電力・計算資源の制約から，暗号化や認証などのセキュリティ機能が十分に実装されない場合がある \cite{Dritsas2025IoTSecurity}．とりわけ，送信フレームに完全性保護（MAC 等）や新鮮性（freshness）を担保する情報（カウンタ／シーケンス番号等）が含まれず，受信側も受理済み履歴を管理しない設計では，攻撃者が傍受した正規フレームを再送するだけで，同一コマンドを任意のタイミングで再実行させうる（リプレイ攻撃）．受信側が「いま届いたフレームが新しい操作か，過去フレームの再送か」を識別できない場合，意図しない制御が繰り返し実行され，安全性や運用上のリスクにつながる．例えば無線リモコン玩具カーでは，操作者が操作していないにもかかわらず走行や灯光が再実行され，周囲の人や物に接触するおそれがある．

自動車の Remote Keyless Entry（RKE），Passive Keyless Entry and Start（PKES）等を対象とした先行研究では，無線信号の傍受・解析・再送や，ジャミング／中継（relay）等と組み合わせた攻撃により，不正な解錠・始動が可能となる事例が報告されている \cite{Francillon2011PKES,Lin2019IndustrialRemote,Garcia2016RKE}．また，安価な SDR とオープンソース解析ツール（Universal Radio Hacker：URH 等）を用いた，無線プロトコルの解析・再送（リプレイ）の実証も報告されている \cite{Anthi2024RFVuln,Mykhaylova2024RCReplay}．

本研究が対象とする 2.4\,GHz 帯無線リモコン玩具カーは，このような低コスト無線リモコンシステムの一例であり，制約下でのリプレイ耐性設計を検討する対象として有用である．本研究では，まず対象玩具カーの無線信号のビット列構造を解析し，ソフトウェア無線によるリプレイ攻撃を通じて，実際に無線リモコン信号のみで車体を制御可能であることを確認する．その上で，後続章において，確率的チャネル環境下で動作するアンチリプレイ機能の設計と評価へと展開する．

\subsection{研究の問題設定}

IPsec などの広域ネットワーク向けセキュリティプロトコルでは，シーケンス番号とスライディングウィンドウ（ビットマップ）によるアンチリプレイ機能が標準化されている \cite{RFC9827,Awaad2014IPsecWindow}．これらの研究では，パケット損失や遅延・乱序が大きい環境において，ウィンドウサイズやシーケンス番号の扱いが可用性に与える影響が議論されているが，対象は主に IP ネットワーク上の双方向通信であり，本研究が対象とするような一方向の低レート無線リモコンシステムとは前提条件が大きく異なる．また，IoT デバイス向けの軽量認証プロトコルやチャレンジレスポンス方式に関する研究も数多く存在するが，これらはセンサノードやゲートウェイ間の双方向通信を前提とするものが多く，送信側からの一方向ブロードキャストを想定した簡易リモコンへの直接適用は難しい．
双方向通信が可能なシステムであれば，チャレンジレスポンス方式により強力なリプレイ防御を実現できるが，片方向の簡易リモコンシステムにそのまま適用することは難しい．

本研究が対象とする 2.4\,GHz 帯無線リモコン玩具カーのような低コスト一方向無線制御システムでは，(1) マイコン資源や電力が限られていること，(2) 無線チャネルにパケット損失や乱序が存在すること，(3) 攻撃者がソフトウェア無線を用いて容易にリプレイ攻撃を実行できること，という制約と脅威が同時に存在する．
ここで「確率的チャネル環境」とは，各フレームが確率 $p_{\mathrm{loss}}$ で消失し，かつ確率 $p_{\mathrm{reorder}}$ で到着順序が乱れる（隣接フレームの入れ替え，あるいは小さなランダム遅延の付与により乱序が生じる）ような簡略化チャネルモデルを指す（定義と実装の詳細は 5.3 節）．

したがって，本研究で解決すべき中心的な問題は，「このような制約付き一方向無線制御システムにおいて，単純なカウンタ方式，スライディングウィンドウ方式，チャレンジレスポンス方式などの防御メカニズムが，上記の確率的チャネルモデル下でどの程度リプレイ攻撃を抑止しつつ，正当な制御コマンドの可用性を維持できるかを明らかにすること」である．そのため，本研究ではモンテカルロシミュレーションにより多数の試行を行い，攻撃成功率と正規受理率を推定し，各防御方式のセキュリティと可用性のトレードオフを評価する \cite{Serena2025SimCyber}．

\subsection{本研究の目的}

本研究の目的は，市販 2.4\,GHz 帯無線リモコン玩具カーをケーススタディとして，低コスト一方向無線制御システムにおけるリプレイ攻撃の成立要因を明らかにし，確率的チャネルモデルの下で軽量なアンチリプレイ機能の有効性を定量評価することを通じて，当該クラスのシステムに適用可能な設計指針を得ることである．ここで「対象とするシステム」とは，実験に用いる特定の玩具カー個体そのものではなく，低計算資源・電池駆動・損失/乱序の存在・盗聴/再送の容易さといった制約を共有する低コスト一方向無線リモコンシステム一般を指す．なお，本論文では入手容易で代表性の高い 2.4\,GHz 帯玩具カーを具体例として用いる．

具体的には，以下の三点を目的とする．
\begin{itemize}
\item 市販の 2.4\,GHz 帯無線リモコン玩具カーを対象として，ソフトウェア無線機器および汎用解析ツールを用いて無線フレームのビット列構造を解析し，リプレイ攻撃により正規リモコンを用いずに車体を制御可能であることを実験的に示す．
\item 解析結果に基づき，一方向無線制御システムに共通する抽象的なシステムモデルおよび攻撃モデルを定式化する．さらに，ベースライン（防御なし）に加えて，ローリングカウンタ方式，スライディングウィンドウ方式，および（双方向通信を仮定した）簡易チャレンジレスポンス方式を，統一的な枠組みで定義する．
\item パケット損失率 $p_{\mathrm{loss}}$ と乱序率 $p_{\mathrm{reorder}}$ により特徴づけられる確率的チャネルモデル（詳細は後述）の下で，上記防御方式の安全性（攻撃成功率）と可用性（正当コマンド受理率）をモンテカルロシミュレーションにより系統的に比較し，方式選択およびパラメータ選択（例：ウィンドウ幅 $W$）に関する設計指針を示す．
\end{itemize}

ここで 4 方式を評価対象として選定した理由は次のとおりである．第一に，防御なし（none）は最小実装であり，評価指標に対する基準線として不可欠である．第二に，ローリングカウンタ方式（rolling）は一方向通信のまま実装可能であり，受信側の状態管理が小さく，低コスト制約下で現実的な軽量方式の代表である．第三に，スライディングウィンドウ方式（window）は，損失・乱序が存在する環境で可用性を確保しやすく，ウィンドウ幅 $W$ によるトレードオフを評価できる代表方式であり，標準化されたアンチリプレイ機能（シーケンス番号＋ウィンドウ）とも対応づけて議論できる．第四に，チャレンジレスポンス方式（challenge）は双方向通信を要するため単純リモコンへ直接適用しにくいが，「追加リンクを許容した場合に到達しうる比較基準（上限）」として位置づけることで，一方向方式の利点（低コスト）と限界（防御強度の下限）を相対化して議論できる．

\subsection{本論文の構成}

本論文の構成は以下のとおりである．
第2章では，本研究と関連する無線リモコンシステムおよびリプレイ攻撃，
ならびにその防御手法に関する先行研究を整理し，本研究の位置付けを明らかにする．
第3章では，対象とする 2.4\,GHz 帯無線リモコン玩具カーのハードウェア構成と
通信仕様を説明し，ソフトウェア無線を用いた信号計測・ビット列解析および
実機リプレイ実験の結果について述べる．
第4章では，一方向無線制御システムを抽象化したシステムモデルと攻撃モデルを定式化し，
ベースライン方式（防御なし）に加えて，ローリングカウンタ方式，
スライディングウィンドウ方式，簡易チャレンジレスポンス方式からなる
複数のリプレイ攻撃防御メカニズムを定義する．
第5章では，パケット損失および乱序を考慮した確率的チャネルモデルと，
前章で定義した各防御メカニズムを実装したモンテカルロシミュレータの構成と
実装方針について述べる．
第6章では，チャネル損失率や乱序率，防御パラメータ（ウィンドウ幅など）を変化させた場合の
実験結果を示し，安全性（攻撃成功率）と可用性（正当コマンド受理率）の観点から
各防御メカニズムの性能を比較・考察する．
第7章では，本研究の成果を総括するとともに，今後の課題と発展の方向性について述べる．

なお，本論文で用いる主要な記号を表\ref{tab:notation} に示す．

\begin{table}[tb]
\centering
\caption{本論文で用いる主な記号}
\label{tab:notation}
\begin{tabular}{l p{55mm}}
\hline
記号 & 意味 \\
\hline
$p_{\mathrm{loss}}$ &
無線チャネルにおけるパケット損失率
（各パケットが独立にドロップされる確率） \\[0.5ex]

$p_{\mathrm{reorder}}$ &
無線チャネルにおけるパケット乱序率
（各パケットが送信順と異なる順序で到着する確率） \\[0.5ex]

$N_{\mathrm{runs}}$ &
モンテカルロシミュレーションの試行回数
（独立なチャネル実現の数） \\[0.5ex]

$N_{\mathrm{legit}}$ &
1 試行あたりに送信される正規コマンド（合法フレーム）の個数 \\[0.5ex]

$N_{\mathrm{attack}}$ &
1 試行あたりに攻撃者が実行するリプレイ攻撃試行の個数 \\[0.5ex]

$W$ &
スライディングウィンドウ方式および受理ウィンドウ方式における
ウィンドウ幅（受信側が保持するシーケンス番号範囲の幅） \\[0.5ex]

$P_{\mathrm{legit}}$ &
全試行にわたる正規受理率
（送信された正規コマンドのうち受信側に受理された割合） \\[0.5ex]

$P_{\mathrm{attack}}$ &
全試行にわたる攻撃成功率
（攻撃者が送信したリプレイフレームのうち受信側に受理された割合） \\[0.5ex]

$J$ &
総合スコア（$J = P_{\mathrm{legit}} - P_{\mathrm{attack}}$）\\
\hline
\end{tabular}
\end{table}

%%%%%%%%%%%%%%%%%%%%%%%%%%%%%%%%%%%%%%%%%%%%%%%%
\section{関連技術および研究動向}

\subsection{無線リモコン玩具および IoT デバイスの概要}

従来の 27/40\,MHz 帯アナログリモコンに対し，近年の 2.4\,GHz 帯デジタルリモコンは
周波数ホッピングやデジタル変調を用いることで，多台数同時制御や長距離通信が可能となっている．
市販の 2.4\,GHz 帯リモコンカー向けシステムでは，低コストマイコンと少数の外付け部品のみで
数十\,m 程度の通信距離を実現でき，また周波数ホッピング方式により
多数の送信機が同一空間で同時に利用できる点が特長である\cite{RC2p4GSlides,TamiyaRCGuide}．
ホビー向け RC カーでは，プロポーショナル制御に対応した 2.4\,GHz デジタルプロポが標準となっており，
エントリモデルでも安定した操作性を備える\cite{TamiyaRCGuide,EURORC2025}．

一方，スマートロックやスマート家電，センサノードなどの無線制御 IoT デバイスも急速に普及している．
IoT Analytics による最新の調査では，世界の接続 IoT デバイス数は 2023 年末時点で約 166 億台，
2024 年末には約 188 億台に達すると報告されており\cite{IoTAnalytics2024}，
家庭内では Wi\mbox{-}Fi，Bluetooth Low Energy（BLE），Zigbee など，
主に 2.4\,GHz 帯を用いる無線方式が広く利用されている\cite{ABIWireless2024,BLEMarket2024}．
これらの多くは，電池駆動を前提とした小型・低消費電力・低コストデバイスであり，
実装容易性を優先して，ベンダ固有の簡易プロトコルや専用フレーム形式が採用されることが多い．

このような低価格無線リモコン玩具と IoT デバイスは，いずれも
「一方向または低レートのコマンド伝送」「2.4\,GHz 帯などの ISM 帯利用」
「専用プロトコルによる閉じたエコシステム」といった共通点を持つ．
一方で，コストや計算資源の制約から，
暗号化やメッセージ認証，リプレイ防御といったセキュリティ機能が
十分に実装されていない場合が多いことが指摘されている\cite{Dritsas2025IoTSecurity}．
さらに，RF 層の脆弱性に関する研究では，安価な SDR 機器を用いて
無線信号を盗聴・再送するだけで機器の動作を乗っ取れるリスクが
具体的な事例とともに報告されている\cite{Anthi2024RFVuln}．
本研究で対象とする 2.4\,GHz 帯無線リモコンカーは，このような
低コスト一方向無線制御システムの典型例であると位置づけられる．

\subsection{2.4\,GHz 帯と FSK 変調の概要}

2.4\,GHz ISM（Industrial, Scientific and Medical）帯は，2400～2483.5\,MHz に割り当てられた産業・科学・医療用の共用周波数帯であり，多くの国や地域で，規定された送信電力や占有帯域幅の条件を満たす限り免許不要で短距離無線システムに利用できる．ETSI EN 300 328 などの規格では，この帯域を用いる広帯域データ伝送装置の送信電力や占有帯域幅，隣接チャネル漏洩電力などが詳細に規定されており，2.4\,GHz 帯は「世界共通で利用できるが混雑した共有リソース」として扱われている \cite{EN300328,NXP_Coexistence}．

低コストの玩具カーやホビー用 RC システム，簡易な IoT デバイス向けには，2.4\,GHz ISM 帯で動作する汎用 RF トランシーバ IC が広く利用されている．代表例として Nordic Semiconductor 社の nRF24L01+ は，2.400～2.4835\,GHz の ISM 帯で動作し，GFSK 変調により 250\,kbps～2\,Mbps のデータ伝送を実現することが仕様書で示されている \cite{Nordic_nRF24L01}．この種のモジュールを用いることで，安価なマイコンと少数の外付け部品のみで 2.4\,GHz 帯デジタルリモコンシステムを構成できる．

FSK（Frequency Shift Keying）変調は，デジタルデータの「0」「1」を搬送波の周波数の切り替えとして表現する周波数偏移変調方式である．最も単純な 2 値 FSK では，ビット 0 と 1 にそれぞれ異なる 2 つの周波数を割り当てる．FSK は搬送波の振幅（包絡）をほぼ一定に保ったまま周波数のみを変化させる一定包絡（constant-envelope）の変調方式である．このため，非線形なパワーアンプ（nonlinear power amplifier：PA）を高効率動作点で用いる送信機でも，振幅情報を担う方式に比べて歪みの影響を受けにくく，低コスト・低消費電力の無線機に適していることが知られている \cite{Sklar_DC}．また，リモートキーレスエントリ（Remote Keyless Entry, RKE）などの短距離リモコン用途では，ASK/FSK 系の単純なディジタル変調方式が広く採用されていることが，実装事例を扱ったアプリケーションノートで報告されている \cite{FSK_RemoteAN,RKE_Overview}．

本研究で対象とする 2.4\,GHz 帯リモコンカーについても，第 3 章で述べるスペクトラム解析およびビット列解析の結果から，2 値 FSK 型のデジタル変調によりコマンドビット列を送信していると推定される．

なお，本論文で後述する評価では，変調方式やアナログ無線回路の詳細な物理層特性は扱わず，確率的チャネルモデルにより「パケットが正しく届くか否か」というリンクレベルの挙動のみを抽象化して取り扱う．

\subsection{無線セキュリティとリプレイ攻撃}

無線通信システムに対するセキュリティ要求は，一般に機密性（Confidentiality），
完全性（Integrity），可用性（Availability）からなる
いわゆる「CIA トライアド」として整理される．
IoT デバイスや無線リモコンの多くは，ユーザ認証情報や個人情報といった
高い機密性を要求するデータを扱うというよりも，
遠隔からの制御コマンドが改ざんされず，なりすましなく実行されること，
および正当な操作が途切れずに行えることが重要である．
そのため，本研究で対象とするような低価格無線リモコンにおいては，
機密性よりもメッセージの完全性と送信元の認証，
そしてチャネル障害時にも一定の可用性を維持することが実用上の主要な課題となる．

リプレイ攻撃は，攻撃者が一度観測した正規の通信フレームを記録し，
後の時刻や別の文脈で再送することにより，
受信側に過去の操作を再度受け入れさせる攻撃である．
ここで重要なのは，攻撃者がフレーム内容を解読・改ざんできなくても，
そのまま「コピーして貼り付ける」だけで攻撃が成立する点である．
暗号化された通信であっても，暗号化前後のビット列がそのまま再利用できる設計になっている場合には，
同一暗号文ブロックを再送するだけでリプレイ攻撃が成立しうる．
したがって，リプレイ攻撃はメッセージの内容ではなく，
新鮮性（freshness）や一意性の欠如を突く攻撃であり，
単に暗号化を導入するだけでは防止できない\cite{NISTNonce,SenWSNSurvey}．

特に，本研究が対象とするような一方向の低レート無線リモコンでは，
受信側がシーケンス番号や時刻，過去に受理したフレームの履歴を管理しない実装が一般的である．
このようなシステムでは，受信機から見て
「今届いたフレームが，本当に新しいコマンドなのか，過去のコマンドの再送なのか」を
判別する手掛かりが存在しない．
その結果，攻撃者がソフトウェア無線機器を用いて無線フレームを盗聴し，
同じフレームを再送するだけで，
正規のリモコンを用いたのと同じ制御を好きなタイミングで繰り返し実行できる．
次節では，この問題に対処するために提案されてきた代表的なリプレイ防御方式を整理し，
本研究で評価対象とするメカニズムとの対応関係を明らかにする．

\subsection{リプレイ攻撃に対する代表的防御方式}

前節までに述べたように，リプレイ攻撃は「過去の正規フレームをそのまま再送信するだけ」で
成立する非常に単純な攻撃である．
このため，多くの無線リモコンシステムやネットワークプロトコルでは，
フレームの「新しさ（freshness）」を保証することによりリプレイを検出・拒否する仕組みが導入されている．
代表的な方式としては，
(1) シーケンス番号やカウンタに基づくローリングコード方式，
(2) スライディングウィンドウを用いたシーケンス番号管理方式，
(3) ノンスを用いたチャレンジレスポンス方式，の三つに大別できる．

まず (1) のローリングコード方式は，自動車の Remote Keyless Entry (RKE) や
ガレージドアオープナなどで広く用いられている\cite{Gesteira2025}．
送信側のキーと受信側の車両 ECU が，
同じ初期値で初期化された疑似乱数生成器（PRNG）やカウンタを内部に持ち，
ボタンを押すたびに「次の」コード値をフレームに埋め込んで送信する．
受信側は「まだ未使用のコード」の集合，
あるいは「現在のカウンタ値より十分大きいかどうか」を確認し，
既に使われたコードや過去のカウンタ値を含むフレームを拒否することで，
単純なリプレイ攻撃を防ぐことができる\cite{Gesteira2025}．
この方式は，一方向通信でも利用可能であり，実装も比較的単純である一方で，
ボタンの空打ちなどによるカウンタのずれ（デスインクロナイゼーション）や，
攻撃者によるジャミング＋記録によって有効な未使用コードを貯められてしまう
といった課題も指摘されている．

次に (2) のスライディングウィンドウ方式は，
IPsec の認証ヘッダ (AH) や Encapsulating Security Payload (ESP) において，
アンチリプレイ機能を実現するために採用されている\cite{rfc6479}．
受信側は，これまでに受理したパケットの最大シーケンス番号 $W_T$ とウィンドウサイズ $W$ を保持し，
区間 $[W_T-W+1, W_T]$ を「受理可能なシーケンス番号の範囲」として管理する．
この範囲の各番号についてはビットマップ等で「すでに受信済みか否か」を記録し，
同じ番号のパケットが再度届いた場合はリプレイとして破棄する\cite{rfc6479}．
一方で，$W_T$ より十分小さいシーケンス番号のパケット（ウィンドウの左側より古いもの）は
「古すぎる」とみなして破棄し，
新しいシーケンス番号を受信したときにはウィンドウ全体を右方向へスライドさせる．
この方式により，一定範囲内のパケットの取りこぼしや順序入れ替わりを許容しつつ，
古いパケットの再利用だけを効率よく検出できる．

最後に (3) のチャレンジレスポンス方式は，
プロトコル参加者が毎回新しいノンス（使い捨て値）をやり取りすることで，
メッセージの新しさを保証する典型的な手法である．
NIST の定義では，ノンスは「再利用されない値であり，
リプレイ攻撃を防ぐためにプロトコルのメッセージに含められる」とされており\cite{NISTNonce}，
ランダム値，タイムスタンプ，シーケンス番号などを単独あるいは組み合わせて用いることが推奨されている．
具体的には，受信側がランダムなチャレンジ値を送り，
送信側がその値とコマンド内容を共有鍵や秘密鍵で認証したレスポンスを返す．
受信側がレスポンスを検証し，かつチャレンジが未使用であれば，
そのメッセージは「今この瞬間に生成されたもの」とみなすことができる．
同じチャレンジを用いたメッセージは二度と受理しないため，
過去のレスポンスを単純に再送しても認証に失敗する．

以上のように，いずれの方式も本質的には
「一度使った値（コード／シーケンス番号／ノンス）は二度と受理しない」という方針に基づいているが，
必要とされる状態管理量や計算コスト，再送やパケット損失に対する耐性は大きく異なる．
特に本研究で対象とする 2.4\,GHz 帯の低価格無線リモコンでは，
マイコンのフラッシュ容量や RAM，電池寿命，
さらに通信が一方向のみであるといった制約が厳しい．
そのため，IPsec で想定されるような大きなスライディングウィンドウや，
複雑なチャレンジレスポンスプロトコルをそのまま適用することは現実的でない．
本研究では，上記 3 つの代表的な考え方を，
低コスト無線リモコンでも実装可能な最小限の形
（ローリングカウンタ＋MAC，簡易スライディングウィンドウ，簡易チャレンジレスポンス）に落とし込み，
パケット損失・順序入れ替わりを伴う確率的チャネル環境下での性能を
シミュレーションにより比較評価する．

\subsection{既存研究と本研究の位置づけ}

本節では，無線リモコンや IoT デバイスに対するリプレイ攻撃および防御手法に関する
既存研究を整理し，本研究の位置づけを明確にする．

まず，自動車分野では，Remote Keyless Entry（RKE）システムを対象とした研究が多数報告されている．
Gesteira-Miñarro らは，車載無線通信に対する代表的なサイバー攻撃事例を体系的に整理し，
リプレイ攻撃やロールジャム攻撃が依然として実用上の脅威であることを指摘している\cite{Gesteira2025}．
また，Quintero らは，実車の RKE キーフォブから SDR で取得した信号を用い，
VGG-16 や SVM による機械学習モデルでリプレイ攻撃を検出する手法を提案している\cite{Quintero2023}．
Abu~Al-Haija らは，ResNet50 を用いたディープラーニングモデルにより，
RKE システムに対する本物の信号と偽のリプレイ信号を高精度に分類できることを示している\cite{AbuAlHaija2022}．
これらの研究は，主に「自動車向けの高度な RKE システム」を対象とし，
深層学習に基づくリプレイ攻撃の検出アルゴリズムに焦点を当てている．

次に，家庭用 IoT デバイスに関しては，Lazzaro らによる REPLIOT ツールが代表的である
\cite{Lazzaro2024,LazzaroREPLIOT}．
彼らは，多数の市販 IoT 機器を対象に，
ローカルネットワーク上の通信を自動収集・再送することでリプレイ攻撃の脆弱性を大規模に評価し，
ローカル通信の保護不足により多くの機器がリプレイ攻撃に対して脆弱であることを示した．
REPLIOT はネットワークレベルのトラフィックを対象とし，
個々の無線物理層フレームのビット列や，
安価なホビー向け 2.4\,GHz リモコンのような「非常に単純なプロトコル」の評価までは踏み込んでいない．

一方，Wireless Sensor Network（WSN）や IoT 認証プロトコルの研究においては，
リプレイ攻撃は古くから基本的な脅威として扱われており，
多数のサーベイや対策プロトコルが提案されている．
Sen は WSN 向けセキュリティに関する包括的サーベイの中で，
暗号化・メッセージ認証・カウンタ等を組み合わせることで
リプレイ攻撃を軽減できることを整理している\cite{SenWSNSurvey}．
Nain らは IEEE~802.15.4 トランシーバのセキュア設計を提案し，
データの新鮮性とリプレイ攻撃耐性を考慮した暗号化・認証機能を組み込むことで，
低消費電力無線リンクのセキュリティを高められることを示している\cite{Nain2017}．
また，Feng らは IoT 向け認証方式として，
タイムスタンプではなくチャレンジレスポンスに基づく軽量な認証・鍵共有プロトコルを提案し，
リプレイ攻撃に対する耐性を形式的手法とシミュレーションの両面から分析している\cite{Feng2017}．
これらの研究は，本研究が評価対象とする「カウンタ＋MAC 型」および
「チャレンジレスポンス型」の防御手法の基本思想を与えるものである．

しかし，上記の既存研究には，以下のようなギャップが残されている．

\begin{itemize}
  \item 多くの研究は，自動車用 RKE や ZigBee/IEEE~802.15.4 など，
        標準化されたプロトコルや高機能な鍵管理機能を前提としており，
        安価な 2.4\,GHz 玩具・ホビー向け無線リモコンを対象とした体系的評価は少ない．
  \item 検出系の研究は深層学習など高度な計算資源を前提としていることが多く，
        極めて制約の厳しい低価格マイコンを前提とした
        「どこまで簡易な防御で，どの程度リプレイ攻撃を抑止できるか」という
        観点での比較評価は十分ではない．
  \item シミュレーション研究は存在するものの，多くはネットワークレベルの
        経路制御や鍵配布プロトコルを対象としており，
        「単一の無線リンクにおけるパケット損失・並べ替えを含む確率的チャネル環境下で，
        複数の防御方式をモンテカルロ的に比較する」という枠組みは限定的である．
\end{itemize}

著者の知る限り，低コスト 2.4\,GHz 無線リモコンを対象に，
パケット損失や順序入れ替わりを含む確率的チャネル環境下で
複数の軽量防御方式を系統的に比較した研究は報告されていない．
本研究は，2.4\,GHz 帯の市販玩具カーに用いられている FSK 無線リモコン信号を実際に解析し，
ビット列レベルでリプレイ攻撃を再現可能な評価環境を構築した点に特徴がある．
そのうえで，既存プロトコルで広く用いられているカウンタ＋MAC 型防御や
チャレンジレスポンス方式の考え方を，低コストデバイス向けに単純化して実装し，
確率的チャネルモデルの下で複数方式の性能を体系的に比較する．
このように，本研究は「高度な車載システム」や「標準化 IoT プロトコル」と
「極めて単純な玩具向け 2.4\,GHz 無線システム」の間を埋める位置づけを持ち，
安価な無線リモコンに対する現実的なリプレイ攻撃対策を検討するための
基礎データを提供することを目的とする．

%%%%%%%%%%%%%%%%%%%%%%%%%%%%%%%%%%%%%%%%%%%%%%%%
\section{市販 2.4\,GHz リモコンカーを用いたリプレイ攻撃事例}

\subsection{対象リモコンカーの仕様と機能}

本研究で対象とするのは，中国製の市販 2.4\,GHz 帯リモコンカーである．
取扱説明書には「2.4\,GHz デジタル無線リモコン」とのみ記載されており，
詳細なプロトコル仕様や暗号化方式は公開されていない．
図\ref{fig:target_car}に，対象とする車体および送信側リモコンの外観を示す．

\begin{figure}[t]
  \centering
  \includegraphics[width=0.80\linewidth]{figures/car.pdf}
  \caption{対象とする市販 2.4\,GHz 帯リモコンカー（車体および送信側リモコンの外観）}
  \label{fig:target_car}
\end{figure}

リモコンには，車両の走行を制御する 2 本のスティックと，
灯光の ON/OFF を切り替えるボタンおよびデモ走行を開始するボタンの，
計 4 つの操作子が搭載されている．左右のスティックはそれぞれ左輪・右輪の
前進／後退を独立に制御する．両輪を同方向に駆動することで直進および後退が
実現され，一方の輪のみを駆動することで旋回動作が得られる．デモボタンを押すと，
前進・後退・左右旋回などの一連の動作が自動的に実行される．灯光ボタンは押すたびに
車体前部の LED の点灯／消灯をトグルし，一度点灯すると次に灯光コマンドを受信するまで
点灯状態を維持する．このように，本対象システムは安価な玩具であるものの，複数の自由度を
持つ走行制御と補助機能（灯光・デモ）を備えた，典型的な 2.4\,GHz 帯無線リモコンシステムの
一例といえる．

\subsection{実験装置および測定環境}

実験に用いた装置は，リモコンおよび車体に加えて，
ソフトウェア無線機 HackRF One，2.4\,GHz 帯用アンテナ，
計測用 PC（Linux），および無線解析ツール URH (Universal Radio Hacker) \cite{URH} である．
HackRF One は受信（RX）と送信（TX）の両モードで使用し，
IQ サンプルの取得と重放信号の送信を同一ハードウェアで行った．
図\ref{fig:exp_system}に実験システムの構成を示す．

\begin{figure*}[t]
  \centering
  \begin{tikzpicture}[
      box/.style={draw, rounded corners, align=center,
                  minimum width=35mm, minimum height=14mm},
      databox/.style={draw, rounded corners, align=center,
                  minimum width=20mm, minimum height=10mm, fill=gray!15},
      small/.style={font=\footnotesize},
      tiny/.style={font=\scriptsize},
      >=Latex
    ]

    % Sender (Remote) and Receiver (Car)
    \node[box] (S) {送信側リモコン};
    \node[box, right=55mm of S] (R) {受信側車体};

    % Attacker below the middle
    \path (S) -- (R) coordinate[midway] (mid);
    \node[box, below=28mm of mid] (A) {攻撃者\\(HackRF + URH)\\[1pt]{\scriptsize RX: 傍受 / TX: 重放}};

    % Recorded data node (between capture and replay)
    \node[databox, below=12mm of mid] (D) {記録データ\\{\scriptsize (IQ/波形)}};

    % Legitimate control path (solid line)
    \draw[-Latex, thick]
      (S) -- node[small, above]{正規制御（2.4\,GHz）} (R);

    % Capture path: Sender -> Data (dashed, same style as replay)
    \draw[dashed, thick, -Latex]
      (S.south) to[bend right=15]
      node[small, left, pos=0.4]{傍受・記録}
      (D.west);

    % Data -> Attacker (dashed)
    \draw[dashed, thick, -Latex]
      (D.south) -- (A.north);

    % Replay path: Attacker -> Receiver (dashed, same style)
    \draw[dashed, thick, -Latex]
      (A.east) to[bend right=25]
      node[small, right, pos=0.6]{重放送信（2.4\,GHz）}
      (R.south);

  \end{tikzpicture}
  \caption{実験システム構成：正規制御経路と攻撃者による傍受・重放経路}
  \label{fig:exp_system}
\end{figure*}

測定は大学の実験室内で実施した．室内には複数台のPCや Wi-Fi アクセスポイントなどの電子機器が稼働しており，
2.4\,GHz 帯には一定のバックグラウンドノイズおよび他機器からの電波が存在する．
信号捕捉時には対象信号の SNR を最大化するため，車体を HackRF One のアンテナの直上にほぼ接触する位置に配置し，
その周辺でリモコンを操作した．これにより環境ノイズの影響を抑えつつ，リモコンから車体への無線信号を HackRF で確実に受信できるよう配慮した．

\subsection{信号キャプチャ手順}

まず，取扱説明書に記載された 2.4\,GHz 帯を起点として，
URH を用いて周波数スペクトルを観測しながらチャネル探索を行った．
リモコンの各ボタンを操作しつつ，2.4\,GHz 周辺を掃引し，
スペクトル上に一時的に立ち上がるピークとボタン操作のタイミングが一致する周波数を探した結果，
本リモコンカーは中心周波数約 2.475\,GHz のチャネルを使用していることが分かった．

以降の測定では，HackRF One を 2.475\,GHz を中心周波数とする受信モードに設定し，
サンプリングレート 2\,Msps，帯域幅 2\,MHz の条件で IQ サンプルを取得した．
各種コマンド（灯光，デモ，前進／後退，左右旋回）について，
ボタンあるいはスティック操作を繰り返し行い，
それぞれ数秒間分の信号を URH プロジェクトとして保存した．
実験室環境では他の 2.4\,GHz 無線システムからの干渉も存在するため，
後続の解析では，対象信号の SNR が高く，
かつフレーム全体が連続して観測できている区間を選別して用いた．

\subsection{ビット列解析とフレーム構造の推定}
取得した IQ サンプルに対して URH（Universal Radio Hacker）を用いて変調解析を行い，
本システムが 2 値周波数偏移変調（binary FSK）であることを確認した．
URH は無線プロトコルの復調・編集・差分解析などを統合的に行えるツール群であり，
無線プロトコル解析の研究用途で広く用いられている\cite{PohlNoackWOOT18}．
URH 上では Samples/Symbol=2，Bits/Symbol=1 とし，
中心周波数 2.475\,GHz，サンプリングレート 2\,Msps，帯域幅 2\,MHz，
周波数偏移（URH 表示値で約 101.562\,kHz）に基づいて復調（demodulation）を実施した．
これらのパラメータは，後述するリプレイ実験における送信波形生成時にも同一値を用いた．
復調設定の画面例を図\ref{fig:urh_settings}に示す．

\begin{figure}[t]
  \centering
  \includegraphics[width=0.98\linewidth]{figures/urh.pdf}
  \caption{URH による復調設定例（Modulation=FSK，Samples/Symbol=2，Bits/Symbol=1）}
  \label{fig:urh_settings}
\end{figure}

\subsubsection{フレーム構造}
複数のコマンドフレームを比較した結果，すべてのコマンドで先頭 19 バイトが恒常的に一致し，
差分は末尾 5 バイト（40 bits）に集中した．以降，末尾 5 バイトを制御フィールド
$\mathrm{tail}$ と呼ぶ．本研究で観測した共通ヘッダ（prefix）は次の通りである．

\begin{quote}
\texttt{20aab824caebda25da7020cf6e76b67cde28c7}
\end{quote}
%
\noindent\hspace{1zw}代表的なコマンドと対応する $\mathrm{tail}$ は表\ref{tab:cmd_tail}に示す．

\begin{table}[t]
\centering
\caption{リモコンコマンドと制御フィールド（末尾 5 バイト）の例}
\label{tab:cmd_tail}
\begin{tabular}{r l l}
\hline
No. & 動作 & $\mathrm{tail}$（末尾 5 バイト）\\
\hline
1  & light（灯光 ON/OFF）        & \texttt{0636b64700} \\
2  & demo（デモ）                & \texttt{3646cd4500} \\
3  & left\_back（左輪後退）       & \texttt{24549a2700} \\
4  & left\_forward（左輪前進）    & \texttt{2252d85780} \\
5  & right\_forward（右輪前進）   & \texttt{2757ff1700} \\
6  & right\_back（右輪後退）      & \texttt{26d64d8f00} \\
7  & both\_forward（両輪前進）    & \texttt{2353fb4780} \\
8  & both\_back（両輪後退）       & \texttt{24d40baf00} \\
9  & LB+RF（逆時針旋回）          & \texttt{2555b93700} \\
10 & RB+LF（順時針旋回）          & \texttt{22d249df80} \\
\hline
\end{tabular}
\end{table}

\subsubsection{XOR 差分（GF(2)）に基づく規則性}
$\mathrm{tail}$（40 bits）を GF(2) 上のビットベクトルとみなすと，
ビット単位 XOR（$\oplus$）は GF(2) の加算に相当する\cite{LinCostello2004}．
本研究では，コマンド間の差分 $\Delta = \mathrm{tail}_a \oplus \mathrm{tail}_b$ を観測することで，
「動作属性（片輪/両輪，前進/後退）を反転させる差分マスク」が安定して存在することを確認した．

具体的には，
左輪方向反転マスク $M_L$，右輪方向反転マスク $M_R$，
両輪同時反転（前進$\leftrightarrow$後退）マスク $M_{FB}$ を次式で得る．
{\small
\begin{align}
M_L   &= \texttt{24549a2700} \oplus \texttt{2252d85780} = \texttt{0606427080}, \notag\\
M_R   &= \texttt{2757ff1700} \oplus \texttt{26d64d8f00} = \texttt{0181b29800}, \notag\\
M_{FB}&= \texttt{2353fb4780} \oplus \texttt{24d40baf00} = \texttt{0787f0e880}. \notag
\end{align}
}

これらは独立コマンドではなく，既知の合法 $\mathrm{tail}$ に XOR することで，
対応する属性（左輪/右輪方向，前進$\leftrightarrow$後退）だけを反転させる変換子として機能する．
例えば，両輪前進 $\mathrm{tail}_{BF}=\texttt{2353fb4780}$ に対して
{\footnotesize
\begin{align}
\mathrm{tail}_{BF}\oplus M_L &= \texttt{2555b93700} \quad(\mathrm{LB+RF}),\notag\\
\mathrm{tail}_{BF}\oplus M_R &= \texttt{22d249df80} \quad(\mathrm{RB+LF}),\notag\\
\mathrm{tail}_{BF}\oplus M_{FB} &= \texttt{24d40baf00} \quad(\mathrm{both\_back})\notag
\end{align}
}
が成立し，表\ref{tab:cmd_tail} の観測値と一致する．

一方で，GF(2) 的に構成できる任意の組合せが受理されるとは限らないことも確認した．
例えば，
{\small
\begin{align}
t_0 &= \mathrm{tail}_{LF}\oplus \mathrm{tail}_{RF}\oplus \mathrm{tail}_{BF} \notag\\
&= \texttt{2252d85780}\oplus\texttt{2757ff1700}\oplus\texttt{2353fb4780}
= \texttt{2656dc0700} \notag
\end{align}
}%
は複数の合法コマンドの XOR で表せるが，$\mathrm{tail}=\texttt{2656dc0700}$ を送信しても車体は反応しなかった（無効/NOP）．
従って，本制御フィールドは差分マスクにより強い線形構造を示す一方で，
受信側は「合法コード語集合」の部分集合のみを受理し，それ以外は破棄する可能性が示唆される（観測範囲）．

\subsubsection{灯光ベクトル $v_{\mathrm{light}}$ の導出と合成}
灯光 $\mathrm{tail}_{\mathrm{light}}=\texttt{0636b64700}$ と $t_0$ を用い，
{\footnotesize
\begin{align}
v_{\mathrm{light}}
&=\mathrm{tail}_{\mathrm{light}}\oplus t_0\notag\\
&=\texttt{0636b64700}\oplus\texttt{2656dc0700}=\texttt{20606a4000}\notag
\end{align}
}%
を灯光付与ベクトルとして導出した．
この $v_{\mathrm{light}}$ を既知の合法走行 $\mathrm{tail}$ に XOR することで，
同一フレーム内で「灯光＋走行」を合成できることを実験的に確認した（次節）．

例として，両輪前進に灯光を付与すると，
{\small
\begin{align}
\mathrm{tail}_{\mathrm{BF+light}}
&=\texttt{2353fb4780}\oplus\texttt{20606a4000}
=\texttt{0333910780}\notag
\end{align}
}%
同様に，片輪前進に灯光を付与する場合も
$\mathrm{tail}'=\mathrm{tail}\oplus v_{\mathrm{light}}$ により構成できる．具体的に，
{\small
\begin{align}
\mathrm{tail}_{\mathrm{LF+light}}
&=\texttt{2252d85780}\oplus\texttt{20606a4000}
=\texttt{0232b21780},\notag\\
\mathrm{tail}_{\mathrm{RF+light}}
&=\texttt{2757ff1700}\oplus\texttt{20606a4000}
=\texttt{0737955700}.\notag
\end{align}
}


\subsection{リプレイ攻撃の実施と結果}
\label{sec:replay_result}

取得したビット列が実際に車体の制御に利用できることを確認するため，
HackRF One および URH の送信機能を用いてリプレイ攻撃を行った．
リプレイ攻撃は，正規通信で観測されたメッセージを記録し，
後から同一（または改変した）信号を再送することで動作を再現する攻撃である\cite{rfc3552}．

\subsubsection{リプレイ手順}
URH 上で各コマンドのビット列（フレーム）を指定し，
受信時と同じ中心周波数 2.475\,GHz，サンプリングレート 2\,Msps，
2 値 FSK の復調結果に基づく FSK パラメータ（周波数偏移およびシンボル長）を用いて
送信波形を生成した（設定例は図\ref{fig:urh_settings}）．
生成した波形を HackRF から送信し，リモコンを一切操作しない状態で車体の動作を観察した．

\subsubsection{リプレイ結果（動作再現）}
その結果，前進／後退／旋回／灯光／デモ等の各コマンドについて，
取得したフレームをそのまま再送するだけで車体が同様に動作することを確認した．
すなわち，本対象は「観測した制御フレームの再送」により動作が再現可能であり，
攻撃者が無線区間を盗聴できる環境では重放（リプレイ）攻撃が成立する．

\subsubsection{改変リプレイ（新たな制御パターンの合成）}
さらに，前節で導出した差分マスク（$M_L, M_R, M_{FB}$）および灯光付与ベクトル $v_{\mathrm{light}}$ を用いることで，
単なる再送にとどまらず，既知の合法フレームから新たな制御パターンを構成できることを確認した．

例として，両輪前進 $\mathrm{tail}_{BF}=\texttt{2353fb4780}$ に灯光を付与した制御フィールドは
$\mathrm{tail}_{BF+light}=\texttt{0333910780}$ となるため，
共通ヘッダ（prefix）と連結して次のフレームを構成した．

\begin{quote}
\texttt{20aab824caebda25da7020cf6e76b67cde28c7 0333910780}
\end{quote}
%
\noindent\hspace{1zw}上記フレームを送信したところ，小車が前進しながら灯光が点灯することが確認された．
同様に，片輪のみ前進させつつ灯光を点灯させるコマンド等，
元のリモコン単体では観測されない組合せパターンも合成可能であることが分かった．
このことは，攻撃者が一度制御フレームを取得すれば，
単なるリプレイに加えて「改変リプレイ」により制御の自由度を拡張し得ることを示している．

\subsubsection{到達距離・安定性に関する考察}
リプレイの成功は送受信距離，周辺干渉，アンテナ指向性，送信利得設定等に依存して変動した．
本研究では定性的に「再送のみで動作再現が可能」であることを主結果として示し，
到達距離や成功率の定量評価は，距離を段階的に変化させて試行回数を確保する実験設計により今後の課題とする．



\subsection{本章のまとめ}

本章では，市販の 2.4\,GHz 帯リモコンカーを対象として，
ソフトウェア無線機と公開ツールのみを用いて無線フレームを取得・解析し，
重放および合成コマンドによる制御が可能であることを示した．
具体的には，URH を用いたスペクトル観測により対象チャネル（2.475\,GHz）を特定し，
2 値 FSK 変調パラメータを推定したうえで，
各種コマンドに対応するビット列を抽出した．
抽出したフレームはいずれも共通のヘッダ部と 5 バイトの制御フィールドから構成されており，
制御フィールドは左右輪の駆動方向および灯光 ON/OFF に対応する
機能ベクトルの XOR により表現できることが分かった．

これらのビット列から合成した送信波形を HackRF One から送信することで，
リモコンを用いずに車体を完全に制御できること，
および元のリモコンには存在しない新たな組み合わせコマンドも実現できることを確認した．
この結果は，本研究で扱うリプレイ攻撃の現実性を裏付けるとともに，
第 4 章以降で構築するシステムモデルおよびシミュレータが，
実機システムに近い前提に基づいていることを示すものである．

%%%%%%%%%%%%%%%%%%%%%%%%%%%%%%%%%%%%%%%%%%%%%%%%
\section{リプレイ攻撃防御メカニズムとシステムモデル}
\label{sec:defense-overview}

\begin{figure*}[tb]
  \centering
  \begin{tikzpicture}[
      box/.style={draw, rounded corners, align=center,
                  minimum width=4cm, minimum height=1.2cm},
      small/.style={font=\footnotesize},
      >=Latex
    ]

    % Sender and Receiver
    \node[box] (S) {送信者 (S)\\正規リモコン};
    \node[box, right=6cm of S] (R) {受信者 (R)\\玩具車本体};

    % Attacker below the middle of S and R
    \path (S) -- (R) coordinate[midway] (mid);
    \node[box, below=2.5cm of mid] (A) {攻撃者 (A)\\HackRF + URH (RX/TX)};

    % Probabilistic wireless channel C (正規チャネル only)
    \draw[-Latex, thick]
      (S) -- node[small, above]{確率的無線チャネル $C$（$p_{\mathrm{loss}},\,p_{\mathrm{reorder}}$）}
      (R);

    % Recording arrow: 信道中点 mid -> Attacker (passive eavesdropping)
    \draw[dashed, -Latex]
      (mid) -- node[small, left]{盗聴・記録}
      (A.north);

    % Replay arrow: Attacker -> Receiver (post-play)
    \draw[dashed, -Latex]
      (A.east) to[bend left=20]
      node[small, right]{事後重放 (post-play)}
      (R.south);

  \end{tikzpicture}
  \caption{一方向無線リモコンシステムとリプレイ攻撃者のシステム・脅威モデル}
  \label{fig:system_threat_model}
\end{figure*}

\subsection{問題の再定式化}

第3章で示したように，市販の 2.4\,GHz 帯リモコンカーは，ソフトウェア無線機と公開ツールのみを用いたリプレイ攻撃に対して極めて脆弱であり，正規利用者とは独立に車体を完全制御できることが分かった．
一般にリプレイ攻撃とは，認証やアクセス制御に用いられる通信メッセージを攻撃者が一度記録し，後の時刻に再送することで，本来想定されていない動作や不正アクセスを引き起こす攻撃である\cite{NISTReplay}．

無線リモコンに対する既存研究においても，433\,MHz 帯の簡易リモコンプロトコルや\cite{Mykhaylova2024RCReplay}，自動車のリモートキーレスエントリ（RKE）システム\cite{AbuAlHaija2022}に対して，記録した無線フレームの単純な再送だけでゲートや車両を解錠できるケースが多数報告されている．
これらのシステムは，コストや省電力を優先するあまり，フレームの新鮮さ（freshness）や一度限りの使用を保証する仕組みが十分に導入されておらず，攻撃者が過去の正規通信をそのまま再利用できてしまう点に脆弱性が存在する．

本研究が対象とする市販リモコンカーの無線システムは，送信側リモコンから受信側車体への一方向リンクであり，ACK や暗号化された双方向ハンドシェイクは存在しない．
攻撃者は，第3章で示したように，正規リモコンと同じ周波数チャネル上でフレームを傍受・蓄積し，任意の時刻に再送する能力を持つと仮定する．
一方で，本研究では，暗号鍵の解読や真正な MAC 値の偽造といった高度な攻撃は扱わず，
過去に観測した正規フレームをそのまま再送する単純なリプレイ攻撃を主な対象とする．
実装上は，実機トレース中の特定コマンドのみを選択的に再生する設定も含まれるが，
いずれも「過去に受信したフレームの再送」に限定し，新たな制御ビット列の合成は本論文の評価範囲外とする．

また，実際の 2.4\,GHz 帯無線リンクでは，パケット損失や重複受信，フレームの順序入れ替わりなど，確率的なチャネル劣化が不可避である．
したがって，防御メカニズムの評価においては，「攻撃成功確率」だけでなく，「正規利用時の使い勝手」を表す指標も同時に考慮する必要がある．
本研究では，バイオメトリクス分野などで広く用いられている\cite{FIDO} false accept rate (FAR) および false reject rate (FRR) の考え方にならい，

\begin{itemize}
  \item 攻撃者によるリプレイフレームを防御機能が誤って受理する確率（FAR に相当），
  \item 正規リモコンからのフレームを防御機能が誤って破棄する確率（FRR に相当）
\end{itemize}%
の2つを主要な評価指標として定義する．
なお，第5章以降では実装上の指標として，
「正規受理率」と「攻撃成功率」を用いる．
前者の $1 - \text{正規受理率}$ が上記の FRR に相当し，
後者が FAR に相当する．
第5章のシミュレーションでは，これらの指標を確率的チャネルモデルの下で定量的に比較し，各種防御方式のトレードオフを明らかにする．

以上を踏まえ，本章では，第3章で扱った具体的な玩具リモコンカーの事例を抽象化し，送信者・受信者・攻撃者・チャネルから成る一般的なシステムモデルとして再定式化する．
まず 4.2 節でチャネル損失やリプレイ攻撃を含む状態遷移モデルを導入し，続いてベースライン方式および複数の軽量防御メカニズム（ローリングカウンタ方式，スライディングウィンドウ方式，チャレンジレスポンス方式など）を，同一の枠組みの中で記述・比較できるよう形式化する．

\subsection{システムモデルとチャネル仮定}

本節では，実機の 2.4\,GHz 玩具車リモコンシステムと，第 5 章で述べる
モンテカルロシミュレータの両方に共通して用いるシステムモデルおよび
チャネル仮定を整理する．
図~\ref{fig:system_threat_model} に，本研究で扱う一方向無線リモコンシステムとリプレイ攻撃者のシステム・脅威モデルを示す．
送信側リモコンを送信者 $S$，
玩具車本体を受信者 $R$，無線信号を傍受・再送する攻撃者を $A$ とする．
$S$ は離散時間インデックス $t = 1,2,\ldots$ に沿って制御フレーム $f_t$ を送信し，
$R$ は受信に成功したフレームのみを到着順に処理する．
各フレーム $f_t$ は，第 3 章で示したように，ライト点灯・左右輪の前進／後退などの
コマンド種別を表す制御フィールドと，防御メカニズムに応じたメタデータ
（ローリングカウンタ値，スライディングウィンドウのインデックス，
チャレンジノンスおよび MAC など）から構成される．
本章では，
ビット列の詳細構造そのものではなく，「1 回の正規コマンド送信につき
長さ固定のフレームが 1 個生成される」という抽象化に留める．

無線チャネルは短距離の 2.4\,GHz ISM 帯リンクであり，多重アクセス，干渉，
フェージングなどによりパケット損失や到着時間のばらつきが生じる．
インターネットや無線ネットワークのパケット損失過程は，時間相関を持つ
バースト損失であることが知られており，拡張ギルバートモデルや
Gilbert--Elliott モデルが広く用いられている\cite{Jiang2000Modeling,Hasslinger2008GE}．
しかし本研究では，玩具車という単一ホップ・低遅延のシナリオを対象とし，
解析の見通しの良さとパラメータ数の少なさを優先して，
各フレームの損失を独立なベルヌーイ試行とみなす
I.I.D. パケット損失モデルを採用する．
具体的には，正規フレームおよび
攻撃フレームのそれぞれについて，損失確率 $p_{\mathrm{loss}}$ に従い独立に
「到達／消失」を決定する．

パケットの順序入れ替えについても，同様に単純化したモデルを用いる．
IP Performance Metrics (IPPM) では，パケット再順序化を測定するためのメトリックが
標準化されており，ルータのキューイングや複数経路伝送によって様々な形の
再順序化が生じることが報告されている\cite{rfc4737}．
本研究では，
その影響をパラメータ $p_{\mathrm{reorder}}$ でモデル化し，送信されたフレーム列に対して
確率 $p_{\mathrm{reorder}}$ で隣接フレームの入れ替え（あるいはランダムな小さな遅延付与）
を行うことで，到着順序の乱れを近似的に再現する．

攻撃者 $A$ は，送信者 $S$ の近傍に受信機を設置して盗聴すると仮定するが，
実験結果が示すように，攻撃側の受信も感度限界や干渉の影響を受ける．
これを反映するため，攻撃者用に独立した損失確率 $p_{\mathrm{attacker\_loss}}$ を導入し，
$S$ が送信したフレームが $A$ に記録されない確率としてモデル化する．
$A$ は記録に成功したフレームのみを保存し，攻撃フェーズにおいて
シミュレータが指定するスケジューリング戦略（事後一括再生や
インライン混入など）に従って再送信する．
再送信フレームも，
正規フレームと同じ $p_{\mathrm{loss}}$ および $p_{\mathrm{reorder}}$ を持つチャネルを通過すると仮定する．

以上のように，本研究のシステムモデルは
(1) 3 つの主体 $S$・$R$・$A$，
(2) i.i.d. パケット損失モデルと単純な順序入れ替えモデル，
(3) 攻撃者固有の記録失敗確率 $p_{\mathrm{attacker\_loss}}$
から構成される．
バースト損失やチャネル状態の遷移を持つ
より現実的なモデル（Gilbert--Elliott チャネルなど\cite{Hasslinger2008GE,Jiang2000Modeling}）は，
評価パラメータの増加と計算コストの観点から本論文では対象外とし，
今後の拡張課題とする．

\subsection{ベースライン方式（リプレイ防御なし）}

本節では，本研究で用いる比較対象として，リプレイ攻撃に対する防御を一切行わないベースライン方式（no\_def モード）を定義する．
これは，第 3 章で解析した玩具小車の実機プロトコルに対応しており，受信側が過去の通信履歴や状態を保持せず，「見た目が正しそうなフレームであれば常に受理する」という単純な方式である．

\subsubsection{受信処理の定義}

ベースライン方式において受信機 $R$ は，チャネルから復調されたフレーム $f$ を受信したとき，次のようなフォーマット検査のみを行う．

\begin{itemize}
  \item 先頭のプレアンブルおよび固定ヘッダが，既知の値と一致しているかどうか
  \item フレーム全長が想定されるバイト数と一致しているかどうか
\end{itemize}

CRC などの簡易なフレーム整合性チェックは行うものの，フレーム中のペイロードが「新しいコマンドかどうか」を判断するためのカウンタやノンスは一切保持しない．
したがって，攻撃者が過去の正規通信を記録し，まったく同じビット列を再送した場合でも，チャネル誤りがなければ正規コマンドと区別されずに受理される．

シミュレーション上は，第 4.2 節で定義した確率的チャネルモデルを用い，ビット誤り・CRC エラーなどによる破棄はすべて「パケット損失確率 $p_{\mathrm{loss}}$」としてまとめて扱う．
そのため，ベースライン方式の受信判定は「フレームがチャネルを通過して受信に成功したかどうか」のみで近似される．
正規送信者からのコマンドフレームと攻撃者からのリプレイフレームが同じスロットで同時に到着した場合には，チャネルのランダム性を反映するため，シミュレータでは一定の確率（本研究では $1/2$）でどちらか一方だけが受理されるものとした．

\subsubsection{評価指標における位置づけ}

ベースライン方式は，リプレイ防御機能を持たない現実の低価格無線システムを近似するものであり，以下の理由から比較対象として採用した．

\begin{itemize}
  \item 玩具小車をはじめとする多数の 2.4\,GHz 帯低価格リモコンでは，コマンドの正当性検査は事実上「ビットパターンが一致するかどうか」と簡易なフォーマット検査に限られており，リプレイ攻撃に対しては無防備であると考えられる．
  \item センサネットワーク向けリンク層フレームワーク TinySec \cite{TinySec} は，設計者自身が「リプレイ保護はリンク層ではなくアプリケーション層が担うべきである」と主張しており，リンク層でリプレイ防御を実装しない構成も許容している．
  \item FlexiSec \cite{JinwalaFlexiSec} のように，リプレイ保護の有無をアプリケーション要求に応じて切り替える構成が提案されていることからも，「リプレイ防御なし」状態を明示的な運用モードとして扱うことには現実的な意味がある．
\end{itemize}

以上より，本論文では，ベースライン方式を「チャネル条件のみを考慮し，過去の受信履歴に基づくいかなる検査も行わない受信機」として定式化し，第 5 章以降で提案する Rolling Counter 方式，Sliding Window 方式，Challenge--Response 方式と比較する際の基準とする．

\subsection{ローリングカウンタ + MAC 方式}

本研究の第3章で解析した玩具小車プロトコルでは，制御フレームの末尾に「ほぼ常に異なる」ビット列が付与されており，これは簡易的なローリングカウンタあるいはシーケンス番号として解釈できる．
この観察に基づき，本節では「ローリングカウンタとメッセージ認証コード（MAC）を組み合わせた防御方式」を，シミュレータ上の抽象モデルとして定義する．
同様の設計思想は，無線センサネットワーク向けの SPINS/SNEP プロトコル\cite{PerrigSPINS} や HMAC に基づく認証付き暗号\cite{RFC2104HMAC}，自動車のリモートキーレスエントリ（RKE）におけるローリングコード方式\cite{Garcia2016RKE} などで広く用いられている．

本方式では，送信側デバイス $T$ と受信側デバイス $R$ が，以下の状態を共有していると仮定する．

\begin{itemize}
 \item 共有秘密鍵 $K$（MAC 計算用）
 \item 送信側カウンタ $ctr_T$ と受信側カウンタ $ctr_R$（初期値は 0）
\end{itemize}

送信側は，各コマンド送信ごとに $ctr_T$ を 1 だけインクリメントし，カウンタ値を含むフレーム全体に対して MAC を計算する．

\[
\begin{aligned}
\mathrm{tag} &= \mathrm{MAC}_K(\mathit{cmd} \,\|\, \mathit{payload} \,\|\, ctr_T),\\
\mathrm{frame} &= (\mathit{cmd},\, \mathit{payload},\, ctr_T,\, \mathrm{tag}).
\end{aligned}
\]

受信側はフレームを受信すると，以下の順で検証を行う．

\begin{enumerate}
 \item まず MAC を検証し，改ざんされていないことを確認する．
 \item MAC が正しければ，カウンタ値 $ctr$ と保持している $ctr_R$ を比較し，
       \[
          ctr > ctr_R
       \]
       のときのみフレームを受理し，状態を $ctr_R \leftarrow ctr$ に更新する．
 \item $ctr \le ctr_R$ の場合はリプレイとみなして破棄する．
\end{enumerate}

この単純なルールにより，攻撃者が過去のフレームを記録して再送信しても，カウンタが「過去の値」である限り，受信側はそれを必ず拒否できる．
一方で，カウンタ値が十分大きくなった場合（オーバーフロー），あるいは何らかの理由で送信側・受信側のカウンタが非同期になった場合には，追加の設計配慮が必要になる．
SPINS に含まれる SNEP では，メッセージごとに増加するカウンタに基づいてリプレイ攻撃を防ぐことを前提としており，カウンタ値が再利用されないことが重要である\cite{PerrigSPINS}．
実運用においてカウンタが尽きる，あるいは大きくずれた場合には，再鍵交換やカウンタ同期のための再送プロトコルなど何らかの仕組みを導入することが必要になると考えられる．
本研究のシミュレータでは，カウンタは 32 bit の符号なし整数として扱い，実験中にオーバーフローしない前提で評価を行った．

ローリングカウンタ + MAC 方式は，理想的なモデルでは「最新カウンタ値より小さいフレームはすべて拒否」するため，リプレイ攻撃に対して非常に強い．
一方で，実際の無線環境ではパケット損失や順序入れ替わりが生じるため，カウンタ値が常に連続して受信されるとは限らない．
DTLS 1.2 では，シーケンス番号とスライディングウィンドウを用いたアンチリプレイ機能が規定されており\cite{RFC6347DTLS}，一定幅のウィンドウの中で「未使用のシーケンス番号のみ受理し，過去に受信した番号を拒否する」ことで，損失耐性とリプレイ耐性のバランスを取っている．
本シミュレータにおけるローリングカウンタ + MAC 方式は，このような実プロトコルを単純化した「ウィンドウ幅 1」のモデルと位置づけられ，理想的な条件下でのリプレイ耐性の上限性能を評価するためのベースラインとして利用する．

\subsection{スライディングウィンドウ方式}

前節のローリングカウンタ + MAC 方式は，常に「最新のカウンタ値より小さい値をすべて拒否する」ことで，高いリプレイ耐性を実現する一方，パケット損失や順序入れ替わりに対して脆弱であるという欠点を持つ．
実際のプロトコルでは，Datagram Transport Layer Security (DTLS)\cite{RFC6347DTLS} や IPsec ESP\cite{rfc4303} において，シーケンス番号とスライディングウィンドウを用いたアンチリプレイ機能が採用されており，一定範囲内の過去番号を「まだ未使用であれば受理する」ことで，損失耐性とリプレイ耐性のバランスを取っている．
本節では，これらの方式を簡略化したスライディングウィンドウ型アンチリプレイ機能を，本研究のシミュレータ向けに抽象化して定義する．

本方式でも，送信側 $T$ は前節と同様に単調増加するシーケンス番号（カウンタ）を各フレームに付与し，MAC によって改ざん検出を行うものとする．
受信側 $R$ は，状態として以下の 2 つを保持する．

\begin{itemize}
 \item これまでに受理した最大シーケンス番号 $S_{\max}$
 \item 長さ $W$ のビットマップ $\mathit{bitmap}$（ウィンドウ幅 $W$）
\end{itemize}

ビットマップ $\mathit{bitmap}[i]$ は，「$S_{\max} - i$ に対応するシーケンス番号のフレームを既に受理したかどうか」を表すフラグである（0: 未受理，1: 受理済み）．
$W$ はウィンドウ幅であり，DTLS 1.2 では 32 以上をサポートし，64 をデフォルトとすることが推奨されている\cite{RFC6347DTLS}．

受信側は，MAC 検証に成功したフレームのシーケンス番号 $S$ を受け取ったとき，次の手順でアンチリプレイ判定を行う\cite{RFC6347DTLS,rfc4303}．

\begin{enumerate}
 \item もし $S \leq S_{\max} - W$ であれば，ウィンドウの左端より古い番号であり，必ずリプレイとみなして破棄する．
 \item もし $S_{\max} - W < S \leq S_{\max}$ であれば，ウィンドウ内の番号である．
   \begin{itemize}
     \item 対応するインデックス $i = S_{\max} - S$ について $\mathit{bitmap}[i] = 1$ ならば，その番号のフレームは既に受理済みであり，リプレイとみなして破棄する．
     \item $\mathit{bitmap}[i] = 0$ ならば，初めて受信した番号であるため，正規フレームとして受理し，$\mathit{bitmap}[i] \leftarrow 1$ と設定する．
   \end{itemize}
 \item もし $S > S_{\max}$ であれば，新しい番号がウィンドウの右側に現れたことを意味する．この場合はフレームを受理し，ウィンドウを右方向へ $(S - S_{\max})$ だけシフトさせる．具体的には，
   \begin{itemize}
     \item ビットマップを右へ $(S - S_{\max})$ ビットシフトし，新たに開いた若い側のビットを 0 で埋める．
     \item $S_{\max} \leftarrow S$ と更新し，$\mathit{bitmap}[0] \leftarrow 1$ と設定する（最新番号は受理済み）．
   \end{itemize}
\end{enumerate}

このアルゴリズムにより，シーケンス番号がウィンドウ幅 $W$ 以内の範囲で多少前後して到着しても，まだ受理されていない番号であれば受理できる一方，一度受理した番号の再送や，十分古い番号の再送は確実に拒否できる．
DTLS 1.2 の仕様では，このようなビットマップ付きスライディングウィンドウにより，重複レコードの検出および破棄が定義されている\cite{RFC6347DTLS}．
本研究のシミュレータでは，玩具リモコンという低レート・短メッセージ長のシナリオを想定し，
$W$ を 1～20 程度の比較的小さな値（具体的には $W \in \{1, 3, 5, 7, 9, 15, 20\}$）に変化させて性能比較を行う．

本研究のシミュレーション実装では，ローリングカウンタ + MAC 方式と同様に，カウンタオーバーフローは評価範囲外とし，32 bit 符号なし整数として扱う．
IPsec ESP では，受信側がシーケンス番号に対して幅 $W$ のスライディングウィンドウとビットマップを保持する実装が想定されている\cite{rfc4303}．
したがって，$W$ を大きくすると受信側で保持すべき状態量も増加し，低リソースなデバイスではメモリ使用量とのトレードオフが生じる．
また，ウィンドウ幅を極端に大きくすると，攻撃者に許容される「先行送信」（極端に大きな番号を先に送りつけてウィンドウを一気に進める攻撃）の余地も広がるため，$W$ の選択は重要な設計パラメータとなる．
本論文では，$W$ を数種類に変化させ，リプレイ攻撃に対する耐性（FAR の低減）と正規利用時の誤拒否（FRR の増加）とのトレードオフを第 5 章で評価する．

\subsection{チャレンジレスポンス方式（ノンスベース認証）}

\subsubsection{基本原理}

チャレンジレスポンス方式は，受信側がランダムなチャレンジ値（ノンス）を生成し，それに対する正しい応答のみを受理することで，メッセージの「新鮮さ（freshness）」を保証する古典的な認証手法である．
ISO/IEC 9798-2 は，共通鍵暗号を用いたエンティティ認証メカニズムを規定しており，タイムスタンプやシーケンス番号，乱数などの時変パラメータを用いて，過去の認証情報が再使用されることを防ぐことを明示している\cite{ISO9798-2}．
乱数ノンスを用いるチャレンジレスポンスは，その代表例である．

典型的なノンスベース・チャレンジレスポンスの手順は以下のとおりである．

\begin{enumerate}
  \item 受信側 $R$ が十分なビット長（例：64ビット以上）のランダムノンス $N$ を生成し，送信側 $S$ に送る（チャレンジフレーム）．
  \item 送信側はコマンド $C$ とノンス $N$，共有鍵 $K$ に基づき，
        \[
           T = \mathrm{HMAC}_K(C \,\|\, N)
        \]
        を計算し，$(C, N, T)$ を応答フレームとして送信する．
  \item 受信側は，同じ鍵 $K$ で $T$ を再計算し，一致すれば
        「$C$ は自分が直前に発行したチャレンジ $N$ に対する正当な応答である」と判定してコマンドを実行する．
        ノンスが一致しない場合や，既に使用済みのノンスに対する応答はすべて破棄する．
\end{enumerate}

このとき，攻撃者が過去に観測した $(C, N, T)$ を再送しても，受信側はすでに異なるノンス $N'$ に切り替えているため，古いノンスに対する応答として拒否できる．
ノンスを一様乱数として生成し，一度使用したノンスを再利用しない限り，攻撃者は\textbf{リプレイ攻撃によって新しい正当応答を生成することはできない}．
McGrew による IV/ノンス生成ガイドライン\cite{mcgrew-iv-gen-03} でも，ノンスを「一度だけ用いる値」として扱うことが再利用攻撃の防止に重要であることが強調されている．
また，IoT 向け認証方式のサーベイ\cite{ElHajj2019IoTAuth} でも，多数の提案方式が乱数ノンスに基づくチャレンジレスポンス構造を採用し，リプレイ攻撃耐性を実現していることが報告されている．

\subsubsection{本研究システムへの適用モデル}

本研究では，低価格 2.4\,GHz 玩具小車のような\textbf{単方向リモコン}ではなく，将来的に双方向リンクを備えた IoT デバイスを想定した\textbf{理想的防御メカニズム}として，チャレンジレスポンス方式をシミュレータ上にモデル化する．
実機の玩具小車では車体からリモコンへの上りリンクが存在しないため，本方式は直接適用できないが，双方向通信が可能になった場合に達成し得るリプレイ耐性の上限を評価するベンチマークとして位置付ける．

本シミュレータにおけるフレーム構造は概略的に次のように定義する．

\begin{itemize}
  \item チャレンジフレーム（$R \rightarrow S$）：
        \texttt{[CHAL\_ID\ |\ NONCE]}
  \item 応答フレーム（$S \rightarrow R$）：
        \texttt{[CMD\ |\ CHAL\_ID\ |\ NONCE\ |\ MAC]}
\end{itemize}

ここで \texttt{CHAL\_ID} はチャレンジを識別する ID，\texttt{NONCE} は $b_{\mathrm{chal}}$ ビット長の乱数ノンス，\texttt{MAC} は共有鍵 $K$ と $(\texttt{CMD}, \texttt{CHAL\_ID}, \texttt{NONCE})$ に対する HMAC 値である．
受信側 $R$ は，各チャレンジに対して「有効な \texttt{CHAL\_ID} と NONCE の組」と有効期限を状態として保持し，次の条件をすべて満たす場合のみ応答フレームを受理する．

\begin{enumerate}
  \item 受信した \texttt{CHAL\_ID} が，現在有効なチャレンジの ID と一致している．
  \item NONCE が，その \texttt{CHAL\_ID} に対して記録済みの値と完全一致している．
  \item MAC が正しく，ペイロードが改ざんされていない．
  \item チャレンジ発行から有効期限内に応答が到着している．
\end{enumerate}

このモデルのもとでは，攻撃者が過去の応答フレームを任意のタイミングで再送しても，チャレンジ ID／ノンスの不一致，あるいは有効期限切れによって必ず拒否される．
また，攻撃者が新しい正当な応答を偽造するためには MAC を破る必要があり，共有鍵が漏えいしない限り計算量的に現実的ではないと仮定する．

\subsubsection{期待される効果と制約}

チャレンジレスポンス方式を導入した場合，以下のような効果が期待される．

\begin{itemize}
  \item 各コマンドが一度限りのノンスと紐付けられるため，リプレイ攻撃に対して最も強力な防御を提供できる．
  \item カウンタの単調増加やウィンドウ管理に依存せず，チャレンジごとの一回限りの検証で完結するため，カウンタ同期問題の影響を受けにくい．
  \item ノンス長とタイムアウト時間を適切に設定すれば，FAR（攻撃フレームの誤受理率）を非常に低く抑えることができる．
\end{itemize}

一方で，本方式には次のような制約がある．

\begin{itemize}
  \item 受信側から送信側へのチャネルが必須であり，玩具小車のような既存の単方向リモコン製品に適用するにはハードウェア構成の根本的な変更が必要になる．
  \item チャネルが不安定な場合，チャレンジまたは応答のいずれかが損失すると，そのペアは成立しないため，正規コマンドであっても FRR（正規フレームの誤拒否率）が増加する．
  \item 乱数生成器やチャレンジ状態管理，タイムアウト処理など，ローリングカウンタ方式よりも実装が複雑になる．
\end{itemize}

本論文では，上記の理想化チャレンジレスポンス方式を，「双方向 IoT デバイスにおけるリプレイ攻撃防御の理論的上限性能」として位置付け，第 5 章のモンテカルロ実験で他方式との性能比較を行う．

\subsection{本章のまとめ}

本章では，第 3 章で示した玩具小車の実機リプレイ攻撃実験を抽象化し，確率的チャネルモデルとアンチリプレイ機能を組み合わせたシステムモデルを構築した．
まず 4.1 節で，本研究が対象とする問題を「リプレイ攻撃に対する耐性（FAR）と正規利用時の使い勝手（FRR）のトレードオフを，確率的チャネル環境下で定量的に評価すること」として再定式化した．
4.2 節では，送信者・受信者・攻撃者の 3 主体と，i.i.d. パケット損失および簡略化した順序入れ替えモデルから成るチャネル仮定を導入し，第 5 章のモンテカルロシミュレーションで用いる共通の基盤モデルを定義した．

その上で，4.3 節ではリプレイ防御を一切行わないベースライン方式（no\_def）を定義し，実機玩具小車や一部のリンク層セキュリティフレームワークに見られる「フォーマット検査のみを行う設計」に対応するモデルとして位置づけた．
4.4 節では，単調増加カウンタと MAC を組み合わせたローリングカウンタ方式を導入し，カウンタ値の単調性を利用することで強力なリプレイ防御を実現できる一方，パケット損失や順序入れ替わりに対しては脆弱であることを指摘した．
4.5 節では，DTLS や IPsec ESP の実装に倣い，シーケンス番号とビットマップ付きスライディングウィンドウを組み合わせた方式を定式化し，ウィンドウ幅 $W$ をパラメータとして，リプレイ耐性と損失耐性のバランスを調整できるモデルとして整理した．

さらに 4.6 節では，乱数ノンスに基づくチャレンジレスポンス方式を，双方向リンクを前提とした理想的防御メカニズムとして定義した．
この方式は，単方向リモコンである現行の玩具小車には直接適用できないものの，将来的に双方向通信が可能な IoT デバイスにおける「リプレイ攻撃防御の理論的上限性能」を与えるベンチマークとして機能する．

本章で定義した 4 種類の方式（no\_def，ローリングカウンタ + MAC，スライディングウィンドウ，チャレンジレスポンス）は，すべて共通のチャネル仮定と攻撃モデルの上で比較可能なように抽象化されている．
第 5 章では，これらの方式をモンテカルロシミュレーションにより評価し，チャネル条件（損失率・順序入れ替わり率）やパラメータ（ウィンドウ幅 $W$ など）の違いが，FAR／FRR および総合スコアに与える影響を定量的に明らかにする．

%%%%%%%%%%%%%%%%%%%%%%%%%%%%%%%%%%%%%%%%%%%%%%%%
\section{モンテカルロシミュレータの設計と実装}

\subsection{シミュレーションの目的と設計方針}

本研究では，2.4\,GHz 帯無線リモコンシステムを対象とし，チャネル損失やパケット順序入れ替え，攻撃者による再生タイミングなど，複数の確率的要素を含む複雑な挙動を評価する必要がある．これらの条件を実機だけで網羅的に変化させ，十分な試行回数を確保することは，時間・コスト・安全性の観点から現実的でない．そのため，本研究ではモンテカルロシミュレーション\cite{Raychaudhuri2008MC,Law2007Simulation}に基づく専用シミュレータを構築し，同一のチャネル条件と攻撃シナリオの下で複数の防御メカニズムを公平に比較することを目的とする．

モンテカルロシミュレーションは，入力パラメータの確率分布を与え，そこから乱数サンプリングを繰り返すことで出力の分布や統計量を推定する手法であり，工学分野や通信ネットワークの性能評価において標準的に用いられている\cite{Raychaudhuri2008MC,Law2007Simulation}．本研究では，各試行において「正規コマンド列」「チャネル損失・乱序系列」「攻撃者の再生スケジュール」を乱数として生成し，それらを複数の防御モードに対して共通に適用することで，チャネル環境のばらつきによる影響を打ち消し，防御方式間の差分のみを際立たせるよう設計する．

公平性の観点からは，評価結果の違いが「たまたま当たりの良いチャネル系列を引いたかどうか」ではなく，「防御メカニズムそのものの構造」に起因するようにする必要がある．そこで本研究では，各モードに対して同一の乱数列を再利用する共通乱数法による分散削減を行い，少ない試行回数でも統計的に有意な差を観測しやすくしている．また，攻撃者の行動モデル（記録可能なフレーム数，記録失敗確率，リプレイのタイミングなど）も明示的なパラメータとして定義し，防御方式だけでなく攻撃条件も含めた再現性の高い比較を可能とする．

再現性の観点からは，シミュレーションに用いる乱数生成器のシード値を固定し，実験パラメータセットを別途文書化することで，同一の結果を第三者が再現できるようにしている．Law によるシミュレーション研究の標準的な手順\cite{Law2007Simulation}に従い，本研究ではモンテカルロ試行回数 $N_{\mathrm{runs}}$，1 回の試行における正規送信数 $N_{\mathrm{legit}}$ とリプレイ試行数 $N_{\mathrm{attack}}$，チャネルパラメータ（損失率 $p_{\mathrm{loss}}$，順序入れ替え率 $p_{\mathrm{reorder}}$），防御メカニズム固有のパラメータ（ウィンドウ幅 $W$ や MAC 長など）を明示し，実験スクリプトから JSON 形式で集計結果を出力する構成とした．

チャネルモデルについては，まず基礎的なベースラインとして，各パケットが一定確率 $p_{\mathrm{loss}}$ で互いに独立に消失するベルヌーイ型の i.i.d.\ 損失モデルを採用する．前方誤り訂正（FEC）の性能評価に関する Neumann らの研究\cite{Neumann2005FEC}は，平均損失率だけでなく損失の時間相関（バースト性）が FEC の有効性に大きく影響することを示し，i.i.d.\ 損失モデルをより複雑なバースト損失モデルと比較するための基準として用いている．本研究でも，同様に i.i.d.\ 損失モデルを最も単純な基準モデルと位置付けた上で，各防御メカニズムの相対的な性能を比較する．

一方，実際のインターネットや無線 LAN では，誤りが時間的に相関を持つバースト損失が観測されており，これを表現するために Gilbert--Elliott モデルなどの有限状態マルコフチャネルが提案されている．Ha{\ss}linger と Hohlfeld による研究\cite{Hasslinger2008GE}は，IP バックボーンのトラフィックトレースに対して Gilbert--Elliott モデルを当てはめ，単純な 2 状態マルコフモデルでも現実のパケット損失パターンを良好に近似できることを示している．本研究では，卒業研究としてのスコープと計算コストを考慮し，本編では i.i.d.\ 損失と単純な順序入れ替えモデルのみに基づく評価に焦点を絞り，Gilbert--Elliott 型チャネルを用いた評価は今後の発展的課題と位置付ける．

\subsection{シミュレータの全体構成}

第4章で定義したシステムモデルと防御メカニズムを統一的に評価するために，
本研究では，これらを再現性高く比較できるモンテカルロシミュレータを実装した．
Serena らによるサイバーセキュリティ分野のレビュー\cite{Serena2025SimCyber}は，
サイバー攻撃や防御メカニズムの評価においてモデリングとシミュレーションが広く用いられていること，
および既存研究を「アプリケーション領域」「対象とする脅威」「用いられるシミュレーション技法」
「シミュレーションの主目的」という四つの観点から整理できることを示している．
本研究のシミュレータも，このような視点を踏まえ，
(1) 対象とするシステムと脅威モデルを明示した抽象化，
(2) 再現可能な実験構成，
(3) 防御メカニズムやチャネルモデルを入れ替えやすいモジュール分割
を設計方針としている．

具体的には，送信者 $S$，受信者 $R$，確率的チャネル $C$，
リプレイ攻撃者 $A$，およびモンテカルロ試行全体を管理する
実験制御部 $E$ の 5 つのコンポーネントに分割して実装した．
各コンポーネントは，フレーム構造体（コマンド種別，カウンタ値，
ノンス，MAC などのフィールドを持つ抽象的なデータ構造）を介してのみ相互作用する．
内部状態（最後に受理したカウンタ値，期待ノンス，ウィンドウ境界など）は各モジュールに閉じ込め，
他モジュールからはフレームの入出力を通じてしか観測されない．
このようなコンポーネント指向の離散事象シミュレーション構成は，
OMNeT++ や ns-3 といったネットワークシミュレータにおいても採用されており\cite{OMNeTppIntro,ns3Tutorial}，
ノード・チャネル・アプリケーションなどの役割ごとにモデルを分離することで，
モデルの再利用性と検証容易性を高めている．
本研究のシミュレータはこの設計思想を簡略化して取り入れ，
防御メカニズムの追加やチャネルモデルの差し替えを局所的な修正で行えるようにした．

実験制御部 $E$ は，コマンドトレースとシミュレーション設定からなる
\texttt{SimulationConfig} を入力として受け取り，
防御モード（\texttt{no\_def}, \texttt{rolling}, \texttt{window}, \texttt{challenge}）と
モンテカルロ試行回数 $N_{\mathrm{runs}}$ を外側のループとしてシミュレーションを駆動する．
具体的には，$E$ は各モードごとに乱数シード，損失確率 $p_{\mathrm{loss}}$，
順序入れ替え確率 $p_{\mathrm{reorder}}$，
攻撃者側損失確率 $p_{\mathrm{attacker\_loss}}$，
攻撃スケジューリング戦略（\texttt{post} または \texttt{inline}）などの
パラメータを \texttt{SimulationConfig} から読み出し，
その設定の下で $N_{\mathrm{runs}}$ 回の試行を繰り返す．
各試行の開始時には，送信者 $S$，受信者 $R$，チャネル $C$，
攻撃者 $A$ の内部状態を初期化し，トレースファイルから
正規コマンド列（前進・後退・左右旋回・停止など）を読み込む．

単一試行の内部では，送信者 $S$ がトレース中の各コマンドを順に取り出し，
選択された防御モードに応じてフレームを生成する．
\texttt{no\_def} モードではペイロードのみを送信し，
\texttt{rolling} および \texttt{window} モードでは単調増加カウンタと MAC を付与する．
\texttt{challenge} モードでは，受信側が事前に生成したノンスを含む
チャレンジレスポンス形式のフレームを構成する．
生成されたフレームは確率的チャネル $C$ に渡され，
$C$ は第5.3節で述べるモデルに基づき損失および順序入れ替えを適用した上で，
到達したフレームを受信者 $R$ に転送する．

リプレイ攻撃者 $A$ は，送信者 $S$ とチャネル $C$ の間に配置され，
$S$ から出力されるフレームを盗聴する役割を持つ．
$A$ は，攻撃者側損失確率 $p_{\mathrm{attacker\_loss}}$ に従って
記録に成功したフレームのみを内部バッファに蓄積し，
\texttt{post} モードでは正規送信完了後に一括して再生し，
\texttt{inline} モードでは正規フレーム列に対して
一定確率でリアルタイムに再生フレームを挿入する．
再生されたフレームも正規フレームと同じチャネル $C$ を通過し，
同一の $p_{\mathrm{loss}}$ および $p_{\mathrm{reorder}}$ が適用される．

受信者 $R$ は，チャネル $C$ を通じて到着した正規フレームと再生フレームを区別せずに受信し，
選択された防御モードに応じて内部状態（最後に受理したカウンタ値，
受理ウィンドウ，期待ノンス値など）を更新しつつ，
各フレームを受理（accept）または拒否（reject）する．
試行終了時には，$R$ は受理した正規フレーム数と
誤って受理した再生フレーム数を $E$ に返す．
$E$ はこれらの値をモンテカルロ試行全体で集約し，
正規受理率および攻撃成功率の平均値・標準偏差を算出して
JSON ファイルおよび表形式で出力する．
このようにシミュレータの全体構成を明確に分離することで，
脅威モデルの各要素（送信者，チャネル，攻撃者，受信者）が
シミュレーションコード上のどのモジュールに対応しているかを容易に追跡でき，
防御メカニズムの設計や実験結果の解釈を透明にすることができる．

\subsection{チャネルモデル}

本節では，モンテカルロシミュレータ内部で用いるチャネルモデルの具体的な確率分布とパラメータ設定を整理する．本研究では，2.4\,GHz 帯の無線リンクを高精度に再現することよりも，防御メカニズム間の相対的な性能差を明確に比較できることを優先し，チャネルを「パケット消失」と「順序入れ替え」を生成する離散時間の確率モデルとして抽象化した．このようにパケットを「そのまま届く」か「完全に失われる」かの二値として扱う抽象化は，RFC~5170 においてパケット消失チャネル（packet erasure channel）として定義されており\cite{RFC5170}，FEC やストリーミング方式の評価で標準的に用いられている\cite{Neumann2005FEC}．本研究も同様の前提に立ち，ビット誤りそのものではなくパケット単位の消失のみを明示的に扱う．

\subsubsection{パケット損失モデル}

各送信フレーム $f_t$ に対して独立なベルヌーイ試行を行い，確率 $p_{\mathrm{loss}}$ でフレームを消失（受信側には全く届かない）させ，確率 $1-p_{\mathrm{loss}}$ で誤りなく配送されるものと仮定する．これは RFC~5170 で想定されているパケット消失チャネルに相当する\cite{RFC5170}．本研究の実験1では，$p_{\mathrm{loss}} \in \{0, 0.05, 0.10, 0.15, 0.20, 0.25, 0.30\}$ の 7 点を一様刻みで掃引し，低損失から中程度の損失環境までをカバーする．これにより，防御機能が損失率の上昇に対してどの程度まで耐性を維持できるかを観測できる．

Neumann らは，平均損失率が同じであっても，損失が時間的に独立に発生する場合と，バースト状にまとまって発生する場合とでは FEC の性能が大きく異なることを示している\cite{Neumann2005FEC}．一方，Ha{\ss}linger と Hohlfeld は，実際の IP バックボーンで測定したトラフィックトレースに対して Gilbert--Elliott 型の 2 状態マルコフモデルを当てはめ，簡単なマルコフモデルでも現実のパケット損失過程を良好に近似できることを報告している\cite{Hasslinger2008GE}．これらの結果は，現実のネットワークにおいて損失過程がしばしばバースト的かつ時間相関を持ち，単純な i.i.d.\ 損失モデルだけでは最悪ケースの挙動を十分に捉えられない可能性があることを示唆している．

本研究のチャネルモデルでは，卒業研究としてのスコープと計算量を考慮し，各パケットの損失を独立なベルヌーイ試行とみなす i.i.d.\ 損失モデルのみを扱う．したがって，本研究で得られる数値結果は「現実の 2.4\,GHz 無線リンクの精密な再現」を意図したものではなく，\textbf{同一の単純化チャネルモデルの下で，防御機能間の相対的な性能差を比較するための指標}として解釈する必要がある．より現実的なバースト損失チャネル（例えば Gilbert--Elliott モデル）への拡張は，将来の課題と位置付ける．

\subsubsection{パケット乱序モデル}

パケット乱序についても，第4章のシステムモデルと整合するように単純化したモデルを採用した．具体的には，チャネルに投入されたフレーム列に対して，確率 $p_{\mathrm{reorder}}$ で「小さなランダム遅延」を付与する遅延キューを導入する．遅延が付与されたフレームは通常のフレーム数個分だけ後ろにシフトして出力されるため，到着順序がわずかに乱れたフレーム列が生成される\footnote{実装上は最大遅延長を有限に制限しており，極端に古いフレームが突然出現するような挙動は生じない．}．
本研究で用いる $p_{\mathrm{reorder}}$ 自体は，RFC~4737 で定義されるいずれかの乱序メトリクスそのものではなく，「パケット列に含まれる乱序パケットの割合」を表す簡略化された確率パラメータである．したがって，$p_{\mathrm{reorder}}$ の値と RFC~4737 のメトリクス値の間に厳密な 1 対 1 対応があるわけではないが，IPPM で用いられる乱序メトリクス\cite{rfc4737}と同様に，「どの程度の乱序が発生している環境か」を粗くパラメトライズする目的で用いる．本研究では $p_{\mathrm{reorder}} = 0$〜$0.3$ 程度の範囲を掃引している．

RFC~4737 は，ネットワークがパケット順序をどの程度保持しているかを評価するための複数のメトリクスを定義しており\cite{rfc4737}，ルータのキューイングや複数経路伝送などに起因する再順序化が現実のネットワークで継続的に発生し得ることを示している．本研究の乱序モデルは，各フレームに独立に乱序の有無を割り当てるだけであり，現実のような「まとまった乱序バースト」や，キュー長・リンク負荷に起因する複雑な相関構造までは再現していない．したがって，本シミュレーションで得られる結果は，乱序に関する絶対頻度や分布形状を厳密に再現することを意図したものではなく，「一定程度の乱序が存在する環境で，各防御手法のロバスト性にどの程度差が現れるか」を比較するための近似モデルである．

\subsubsection{攻撃者チャネルとの関係}

攻撃者 $A$ が受信・記録するチャネルについては，第4章で導入した記録失敗確率 $p_{\mathrm{attacker\_loss}}$ を用いる．すなわち，送信者 $S$ が送信した各フレームについて，確率 $p_{\mathrm{attacker\_loss}}$ で $A$ はフレームを全く受信できず，確率 $1 - p_{\mathrm{attacker\_loss}}$ で完全なビット列として記録できるものとする．記録に成功したフレームのみが攻撃フェーズで再送信候補となり，再送信フレーム自体は正規フレームと同じ $p_{\mathrm{loss}}$ および $p_{\mathrm{reorder}}$ を持つチャネルを通過すると仮定する．

本研究の主目的は「攻撃者にとって最も有利な状況」で各防御機能の性能を評価することであるため，第5章の実験では特に断りがない限り $p_{\mathrm{attacker\_loss}} = 0$ とし，攻撃者は全ての正規フレームを完全に記録できると仮定する．実際の 2.4\,GHz 環境では攻撃者側の受信も感度限界や干渉の影響を受けるため，多くの場合 $p_{\mathrm{attacker\_loss}} > 0$ であり，本シミュレーションの評価は防御機能にとって保守的（悲観的）な側に振れていると解釈できる．

以上のように，本研究のチャネルモデルは i.i.d.\ 損失と簡易乱序のみを扱う比較的単純なモデルであり，現実の 2.4\,GHz 無線リンクにおける全ての現象（バースト損失，長期相関，マルチパスに起因する複雑な遅延分布など）を再現するものではない．しかし，既存研究が示すように損失分布や乱序の統計的性質がプロトコル性能に与える影響は大きい\cite{Neumann2005FEC,Hasslinger2008GE}ため，本研究ではモデルの前提を明示した上で，\textbf{絶対的な数値そのものではなく，各防御手法間の相対的な傾向を評価指標とする}という立場をとる．

\subsection{出力ログと解析方法}

本シミュレータは，各ラン $r \ (1 \le r \le N_{\mathrm{runs}})$ を独立なチャネル実現とみなし，
ランごとに合法フレームおよびリプレイフレームの受理状況を集計する．
記号の定義で示したとおり，1 ランあたりに送信される合法コマンド（正規フレーム）の個数を
$N_{\mathrm{legit}}$，
攻撃者が送信するリプレイ攻撃試行の個数を
$N_{\mathrm{attack}}$ とする．
さらに，各ラン $r$ において
\begin{itemize}
  \item $A^{(r)}_{\mathrm{legit}}$：ラン $r$ で防御機能によって受理された合法フレーム数
  \item $A^{(r)}_{\mathrm{attack}}$：ラン $r$ で防御機能に誤って受理されたリプレイフレーム数
\end{itemize}
を記録する．
これらのカウンタから，ラン $r$ における
正規受理率（Legitimate Acceptance Rate, LAR）および
攻撃成功率（Attack Success Rate, ASR）を次式で定義する．
\begin{align}
  L^{(r)} &= \frac{A^{(r)}_{\mathrm{legit}}}{N_{\mathrm{legit}}}, \label{eq:def_lar_single_run} \\
  A^{(r)} &= \frac{A^{(r)}_{\mathrm{attack}}}{N_{\mathrm{attack}}}. \label{eq:def_asr_single_run}
\end{align}
ここで $L^{(r)} \in [0,1]$ および $A^{(r)} \in [0,1]$ であり，$A^{(r)}$ はラン $r$ における攻撃成功率（ASR）を表す．
それぞれ「合法トランザクションがどれだけ受理されているか」と
「攻撃トランザクションがどれだけ誤って受理されているか」を表す．
認証分野では，
FIDO 仕様書などで False Accept Rate (FAR) および
False Reject Rate (FRR) が，
生体認証トランザクションに対する誤受理率・誤拒否率として
「全トランザクション数に対する比率」で定義されている\cite{FIDO}．
本研究における ASR は，攻撃者のリプレイ要求を
「偽の本人要求」とみなしたときの FAR に対応し，
$1 - L^{(r)}$ は FRR に対応する量である．
なお，FIDO の FAR/FRR はユーザ認証トランザクション単位の誤受理率・誤拒否率として定義されるのに対し，本研究の LAR/ASR は「フレーム単位の受理・拒否」に対する比率であり，評価の粒度は異なる．
ただし，「受理・拒否という二値判定に対する誤り確率の組」としては同じ構造を持つため，FIDO 仕様で一般的な「安全性と利便性のトレードオフを二つの指標で表す」という考え方を本研究の文脈に応用している．
したがって，
LAR/ASR を用いることで，
FIDO 仕様と同様の枠組みで
「安全性（攻撃成功率の低さ）」と
「利便性（正規利用時の拒否の少なさ）」を同時に評価できる．

次に，各防御方式とパラメータ設定ごとに
$N_{\mathrm{runs}}$ 個のランから得られた
$\{L^{(r)}\}_{r=1}^{N_{\mathrm{runs}}}$ および
$\{A^{(r)}\}_{r=1}^{N_{\mathrm{runs}}}$ を
互いに独立かつ同分布な標本とみなし，
モンテカルロ推定量としての平均値とばらつきを評価する．
具体的には，正規受理率と攻撃成功率の標本平均を
\begin{align}
  \bar{L} &= \frac{1}{N_{\mathrm{runs}}} \sum_{r=1}^{N_{\mathrm{runs}}} L^{(r)}, \label{eq:lar_mean} \\
  \bar{A} &= \frac{1}{N_{\mathrm{runs}}} \sum_{r=1}^{N_{\mathrm{runs}}} A^{(r)} \label{eq:asr_mean}
\end{align}
と定義し，対応する標本分散（不偏分散）を
\begin{align}
  s_L^2 &= \frac{1}{N_{\mathrm{runs}} - 1}
          \sum_{r=1}^{N_{\mathrm{runs}}} \left( L^{(r)} - \bar{L} \right)^2, \\
  s_A^2 &= \frac{1}{N_{\mathrm{runs}} - 1}
          \sum_{r=1}^{N_{\mathrm{runs}}} \left( A^{(r)} - \bar{A} \right)^2
\end{align}
とする．
$N_{\mathrm{runs}}$ が十分大きい場合，
中心極限定理により
$\bar{L}$ および $\bar{A}$ はそれぞれ
ほぼ正規分布に従うことが知られており\cite{Law2015Simulation}，
これを用いると 95\% 信頼区間は近似的に
\begin{align}
  \bar{L} \pm 1.96 \frac{s_L}{\sqrt{N_{\mathrm{runs}}}}, \\
  \bar{A} \pm 1.96 \frac{s_A}{\sqrt{N_{\mathrm{runs}}}}
\end{align}
と表すことができる．
ここで 1.96 は標準正規分布における
95\% 信頼水準に対応する臨界値であり，
離散イベントシミュレーションの出力解析において
広く用いられている\cite{Law2015Simulation}．
本研究では主として
\eqref{eq:lar_mean}, \eqref{eq:asr_mean} による
平均値とそのパラメータ依存性に着目し，
第6章以降の図はいずれも
$\bar{L}$ および $\bar{A}$ の値をプロットして比較を行う．

実装上は，各ラン $r$ ごとに
\begin{itemize}
  \item 防御モード（Rolling Counter, MAC, Sliding Window, Challenge--Response 等），
  \item $p_{\mathrm{loss}}$，$p_{\mathrm{reorder}}$，ウィンドウサイズ $W$，
  \item ラン番号 $r$，
  \item $A^{(r)}_{\mathrm{legit}}$，$A^{(r)}_{\mathrm{attack}}$，
  \item $L^{(r)}$，$A^{(r)}$
\end{itemize}
を 1 行のレコードとして CSV ファイルに出力する．
その後，Python（pandas, matplotlib）を用いた後処理スクリプトにより，
各設定ごとに \eqref{eq:lar_mean}, \eqref{eq:asr_mean} に基づく平均値を計算し，
第6章で示すグラフを生成した．
ログを CSV 形式で保存し，
乱数シードとともに管理することで，
実験条件の変更や再解析を容易にすると同時に，
結果の再現性も確保している．


\subsection{妥当性確認}
\label{subsec:sim-validation}

本研究で用いるモンテカルロシミュレータの出力が信頼できることを確認するために，
Sargent および Kleijnen によるシミュレーションモデルの検証・妥当性確認手法\cite{Sargent2013VVA,Kleijnen1995VV}を参考に，
(1) 実装レベルの検証（verification），(2) 理論的に解析可能な極限ケースとの比較，
(3) 乱数シードおよび試行回数に対する感度解析，という三つの観点から妥当性確認を行った．

まず実装レベルの検証として，送信側・受信側・攻撃者の各モジュールについて，
ごく少数のフレーム列を用いたテストシナリオを作成し，ログを逐次追跡しながら
内部状態が設計どおりに遷移しているかを確認した．
特にスライディングウィンドウ方式については，
IPsec ESP のアンチリプレイ機能\cite{rfc4303}と同様に，
受信側が保持する最大シーケンス番号 $S_{\max}$ と長さ $W$ のビットマップが
常に「過去 $W$ 個の受理済みパケット」を正しく表現していることを，
単調増加・乱序・重複を含む小さなシーケンス列を入力して検証した．
ローリングカウンタ + MAC 方式およびチャレンジレスポンス方式についても，
第4章で定義した状態遷移どおりにカウンタ値や nonce，MAC 検証結果が更新されることを，
単一ランの詳細ログを用いたコード・ウォークスルーによって確認した\cite{Kleijnen1995VV}．

次に，理論的に解析可能な極限ケースとの比較を行った．
チャネル損失・乱序が存在しない $p_{\mathrm{loss}} = 0, p_{\mathrm{reorder}} = 0$ の条件では，
各方式の正規受理率および攻撃成功率は理論的に求めることができる．
例えば無防御方式では，全てのフレームがそのまま受理されるため，
Legitimate Acceptance Rate（正規受理率）および Attack Success Rate（攻撃成功率）は
いずれも 1 になると期待される．
一方，ローリングカウンタ + MAC，スライディングウィンドウ，
チャレンジレスポンス方式では，攻撃者が再送できるのは過去に一度正しく受理された
フレームに限られ，かつ受信側はシーケンス番号や nonce を用いて
「古い／重複した」フレームを拒否するため，
完全なチャネル下では理論的には攻撃成功率 0，正規受理率 1 となる．
本シミュレータにおいて同じ条件でモンテカルロ実験を行ったところ，
得られた各指標はこれらの解析値と 95\% 信頼区間の範囲内で一致しており，
極端なパラメータ設定においても実装が期待どおりに動作していることを確認した．
このように，理論的に振る舞いが分かっている境界条件・極端条件で
シミュレーション結果を検証することは，妥当性確認の代表的な手法である\cite{Sargent2013VVA,Kleijnen1995VV}．

最後に，乱数シードおよび試行回数に対する感度解析を行った．
本研究では各実験条件について $N_{\mathrm{runs}} = 200$ 回のモンテカルロランを実行し，
各ランで $N_{\mathrm{legit}} = 20$ 回の正規送信と
$N_{\mathrm{attack}} = 100$ 回のリプレイ攻撃試行を行う．
各フレームの受理／拒否を互いに独立なベルヌーイ試行と近似すると，
正規受理率の試行総数は $N_{\mathrm{runs}} N_{\mathrm{legit}} = 4000$，
攻撃成功率の試行総数は $N_{\mathrm{runs}} N_{\mathrm{attack}} = 2\times 10^4$ である．
二項分布に基づく最悪の場合（成功確率 $p=0.5$）の標準偏差は，
前者で約 $0.79\%$，後者で約 $0.35\%$ となり，
95\% 信頼区間の幅はそれぞれおおよそ $\pm 1.55\%$，$\pm 0.69\%$ に相当する．
実際に乱数シードを複数パターンに変えて同一条件のシミュレーションを繰り返したところ，
得られた指標値のばらつきは上記の理論的な信頼区間内に収まっており，
採用した試行回数が Monte Carlo 推定量の精度として十分であることを確認した．
このように，シミュレーション結果の統計的誤差の上限を理論的に評価し，
異なる乱数シードやパラメータに対する感度を調べることは，
シミュレーションモデルの妥当性評価において広く用いられている\cite{Sargent2013VVA,Kleijnen1995VV}．

以上の検証結果から，本研究で用いるモンテカルロシミュレータは，
本論文で扱うパラメータ範囲において，4 種類の防御メカニズムの相対的な性能差を
議論するのに十分な精度と一貫性を有していると判断した．

%%%%%%%%%%%%%%%%%%%%%%%%%%%%%%%%%%%%%%%%%%%%%%%%
\section{実験結果および考察}
\label{sec:exp-results}

\subsection{実験条件の整理}
\label{sec:exp-overview}

本章では，4種類の防御メカニズム
（防御なし（none），ローリングカウンタ方式（rolling），
スライディングウィンドウ方式（window），
チャレンジレスポンス方式（challenge））を，
共通のモンテカルロ実験フレームワーク上で比較する．
シミュレータの実装コードおよび代表的な実験パラメータの設定例は，
著者による公開リポジトリに整理しており\cite{ReplaySimReadmeJP}，
第三者による再現実験が可能な形で公開している．

各防御方式に対して，1 つのチャネル条件
$\theta = (p_{\mathrm{loss}}, p_{\mathrm{reorder}}, W, \ldots)$ ごとに
$N_{\mathrm{runs}} = 200$ 回の独立なモンテカルロ試行を行う．
1 回の試行では，正規リモコンから
$N_{\mathrm{legit}} = 20$ フレームを送信し，
攻撃者は高々 $N_{\mathrm{attack}} = 100$ 回のリプレイを行う設定とした．
したがって，各構成あたりの総試行数は，
正規フレームが最大 $200 \times 20 = 4000$ 回，
リプレイ攻撃が最大 $200 \times 100 = 20000$ 回となる．
この規模の試行数であれば，数\%オーダの確率量についても
統計的に安定した推定が得られる．
多数回の試行に基づいて誤受理率・誤拒否率を統計的に評価するという考え方は，
生体認証における FIDO Biometrics Requirements でも採用されている\cite{FIDO-Bio}．
同文書では，False Accept Rate (FAR)，False Reject Rate (FRR) などの指標を，
十分な試行回数に基づいて推定することが求められており，
本研究における実験設計も，このような確率的評価の枠組みに倣っている．

評価指標としては，次の 2 つの確率量を用いる．
\begin{itemize}
  \item \textbf{正規受理率}（Legitimate Acceptance Rate, LAR）：
        正規リモコンから送信されたフレームのうち，
        受信側で受理されたフレームの割合．
        これは $1 - \mathrm{FRR}$ に相当し，
        ユーザビリティ（誤拒否の少なさ）を表す指標である．
  \item \textbf{攻撃成功率}（Attack Success Rate, ASR）：
        攻撃者によって再生されたフレームのうち，
        受信側で誤って受理されたフレームの割合．
        これは $\mathrm{FAR}$ に対応し，
        安全性（誤受理の少なさ）を表す指標である．
\end{itemize}
このように，本研究の LAR／ASR は，
FIDO Biometrics Requirements で用いられる FAR／FRR と
同じ枠組みで解釈できる\cite{FIDO-Bio}．

チャネルモデルについては，
パケット損失率 $p_{\mathrm{loss}}$ とパケット並び替え率 $p_{\mathrm{reorder}}$ を独立に設定し，
\[
  0 \leq p_{\mathrm{loss}} \leq 0.30,\qquad
  0 \leq p_{\mathrm{reorder}} \leq 0.30
\]
の範囲で変化させる．
これは，無線リンクや混雑した IP ネットワークにおいて現実的に観測される
0～数十\% 程度のパケット損失・並び替えを包含することを意図したものである．
リアルタイムメディア向け輻輳制御の評価指針である RFC~8868 では，
エンドツーエンド損失率として 0\%，1\%，5\%，10\%，20\% の条件での評価が
代表例として挙げられており\cite{RFC8868}，
本研究の設定はそれらを含むように上限を 30\% に拡張した形となっている．

本シミュレータのチャネル実装では，
README に明示されているように\cite{ReplaySimReadmeJP}，
パケット損失は確率 $p_{\mathrm{loss}}$ の独立ベルヌーイ過程としてモデル化し，
パケット並べ替えは確率 $p_{\mathrm{reorder}}$ の独立事象としてモデル化する．
すなわち，各フレームの消失・並び替えは時間的に独立であり，
バースト損失や物理層フェージングは明示的には扱わない．
RFC~8868 でも，最低限の評価モデルとして独立ランダム損失の利用が推奨される一方で\cite{RFC8868}，
より現実的なバースト損失のモデルとして
Gilbert--Elliott 型の 2 状態マルコフ損失モデルが挙げられている\cite{Hasslinger2008GE}．
本研究ではまず，解析と実装が単純で再現性の高い独立損失モデルを基準として採用し，
Gilbert--Elliott モデルなどの相関損失モデルを導入した場合との比較は
今後の発展的課題とする．

\subsubsection{理想チャネルにおける実装確認}
\label{subsec:ideal-verify}

防御メカニズムおよび評価指標の実装が仕様どおり動作していることを確認するため，
まずパケット損失率 $p_{\mathrm{loss}} = 0$，
パケット並べ替え率 $p_{\mathrm{reorder}} = 0$ の
理想チャネル条件における挙動を検証した．
この条件ではチャネル側の劣化が存在しないため，可用性とセキュリティの違いは
防御ロジックそのものに起因する．
理想チャネル検証では統計的なばらつきをより小さく抑えるため，
この条件に限り $N_{\mathrm{runs}} = 500$ とした
（結果はシミュレータの出力および README に整理されている\cite{ReplaySimReadmeJP}）．

図\ref{fig:baseline_attack} は，理想チャネルにおける
4 方式（none, rolling, window, challenge）の攻撃成功率（ASR）を比較したものである．
none モードでは正規受理率（LAR）および攻撃成功率（ASR）の両方がほぼ $100\%$ となり，
記録したフレームをそのまま再送するだけで全てのリプレイ攻撃が成功することが分かる．
一方，rolling, window, challenge の 3 方式では，
LAR は $100\%$ を維持しつつ ASR は $0\%$ に抑えられており，
理想チャネルにおいては防御メカニズムが設計どおりに動作していることが確認できる．

このことから，本研究で用いる防御ロジックの実装と，
LAR/ASR の定義および計測手順が妥当であることが検証された．
以降の実験では，$p_{\mathrm{loss}} = 0$ の点を
理想チャネルにおける性能上限（ベースライン）とみなし，
そこからパケット損失や並び替えを導入した場合の性能劣化を評価する．

\begin{figure}[t]
  \centering
  \includegraphics[width=\linewidth]{figures/baseline_attack.pdf}
  \caption{理想チャネル
    $(p_{\mathrm{loss}} = 0,\; p_{\mathrm{reorder}} = 0)$
    における攻撃成功率（ASR）の比較}
  \label{fig:baseline_attack}
\end{figure}

\begin{table}[t]
  \centering
  \caption{実験構成のまとめ（$p_L$: 損失率，$p_R$: 乱序率）}
  \label{tab:exp_overview}
  \footnotesize
  \setlength{\tabcolsep}{3pt}
  \begin{tabular}{c l l}
    \toprule
    実験 & 掃引パラメータ & 固定パラメータ \\
    \midrule
    1    & $p_L \in \{0, 0.05, \dots, 0.30\}$ & $p_R = 0$ \\
    2    & $p_R \in \{0, 0.05, \dots, 0.30\}$ & $p_L = 0.10$ \\
    3    & $W \in \{1,3,5,7,9,15,20\}$ & $p_L = p_R = 0.15$ \\
    Base & なし & $p_L = p_R = 0$, $N_{\mathrm{runs}} = 500$ \\
    \bottomrule
  \end{tabular}
\end{table}

以降の各小節では，表~\ref{tab:exp_overview} に対応する個々の実験結果を示し，
LAR/ASR および総合スコアの観点から防御方式間の差異を考察する．

\subsection{実験1：パケット損失の影響}
\label{sec:exp-loss}

本節では，パケット損失率 $p_{\mathrm{loss}}$ が各防御メカニズムの
可用性とセキュリティに与える影響を評価する．
実験1では，第\ref{sec:exp-overview}節で述べた共通条件のもと，
パケット並び替えを考慮しない $p_{\mathrm{reorder}} = 0$ とし，
\[
  p_{\mathrm{loss}} \in \{0.00, 0.05, 0.10, 0.15, 0.20, 0.25, 0.30\}
\]
の 7 点について，各条件あたり $N_{\mathrm{runs}} = 200$ 回の
モンテカルロ試行を行った．
このとき $p_{\mathrm{loss}} = 0$ は
第\ref{subsec:ideal-verify}節で検証した
理想チャネルベースラインに対応しており，
チャネル劣化が存在しない場合に達成し得る性能の上限を与える．

図\ref{fig:p_loss} に，4 種類の防御モード
（none, rolling, window, challenge）における
正規受理率（LAR，左軸）および攻撃成功率（ASR，右軸）の変化を示す．
いずれも第\ref{sec:exp-overview}節で定義した LAR/ASR を用いており，
FIDO Biometrics Requirements において
生体認証性能の評価指標として用いられている
False Reject Rate (FRR) および False Accept Rate (FAR) と
同じ枠組みで解釈できる\cite{FIDO-Bio}．

\begin{figure}[t]
  \centering
  \includegraphics[width=\linewidth]{figures/p_loss_dual.pdf}
  \caption{パケット損失率と正規受理率（LAR）および攻撃成功率（ASR）の関係
  （$p_{\mathrm{reorder}} = 0$）。
  実線は LAR（左軸），破線は ASR（右軸，対数スケール）を示す。}
  \label{fig:p_loss}
\end{figure}

まず $p_{\mathrm{loss}} = 0$（理想チャネル）の点に着目すると，
4 つの防御モードすべてで LAR はほぼ $100\%$ であり，
防御メカニズム自体が追加の誤拒否を生じていないことが分かる．
一方，none モードでは ASR もほぼ $100\%$ であり，
チャネルに損失や並び替えが全く存在しない状況でも，
防御機能が無ければリプレイ攻撃は常に成功してしまう．
これに対して rolling, window, challenge の 3 方式では，
$p_{\mathrm{loss}} = 0$ において ASR は $0\%$ となっており，
理想チャネルであれば可用性を一切損なうことなく
リプレイ攻撃を完全に遮断できることが確認できる
（第\ref{subsec:ideal-verify}節参照）．

次に，$p_{\mathrm{loss}}$ の増加に伴う LAR の変化を見ると，
いずれの方式でも正規受理率はほぼ単調に低下し，
$p_{\mathrm{loss}} = 0.30$ のときにはおおよそ $70\%$ 前後に収束している
（例として防御なし（none）70.3\%，rolling 70.3\%，window 69.5\%，challenge 70.0\%）．
これは，各フレームが確率 $p_{\mathrm{loss}}$ で独立に失われる
単純なパケット消失モデルを仮定した場合に，
受信成功確率が概ね $1 - p_{\mathrm{loss}}$ に比例して低下するという
直感的な理解と整合的である．
すなわち，実験1の範囲では，可用性低下のほとんどは
チャネル側のパケット損失に起因しており，
防御ロジック固有の追加損失はわずかであると解釈できる．

一方，ASR の挙動を見ると，
$p_{\mathrm{loss}}$ を $0.30$ まで増加させても，
none モードでは ASR が約 69.7\% と依然として高い値を保っているのに対し，
rolling, window, challenge の 3 方式では
$p_{\mathrm{loss}} = 0.30$ においても
ASR がそれぞれ 0.4\%，1.8\%，0.3\% に抑えられている．
すなわち，攻撃者はチャネル上の一部フレームを失うものの，
防御が無い場合には「生き残った」フレームを再送するだけで
高い攻撃成功率を維持できるのに対し，
防御方式ではカウンタやウィンドウ，nonce によって
古いフレームが体系的に拒否されるためである
（第\ref{sec:defense-overview}章参照）．

以上より，実験1からは，
パケット損失は主として可用性（LAR）の低下として現れ，
適切に設計された防御メカニズムであれば，
広い損失率範囲にわたって攻撃成功率（ASR）を低い水準に維持できることが分かる．
理想チャネルに対応する $p_{\mathrm{loss}} = 0$ の点は，
この挙動を理解するための性能上限（ベースライン）として機能している．

\subsection{実験2：パケット並び替えの影響}

\label{sec:exp-reorder}

本節では，固定のパケット損失率のもとでパケット並び替え（乱序）が各防御メカニズムに与える影響を評価する．
実験2では，第\ref{sec:exp-overview}節と同じモンテカルロ条件を用い，
パケット損失率を $p_{\mathrm{loss}} = 0.10$ に固定し，
\[
  p_{\mathrm{reorder}} \in \{0.00, 0.05, 0.10, 0.15, 0.20, 0.25, 0.30\}
\]
の 7 点について，各条件あたり $N_{\mathrm{runs}} = 200$ 回の試行を行った．
ここで $p_{\mathrm{reorder}}$ は，送信順序と異なる順序で到着するフレームの割合を表すパラメータである．

図\ref{fig:p_reorder} に 4 方式における LAR（左軸）および ASR（右軸）の変化を示す．
実験1と同様，評価指標は第\ref{sec:exp-overview}節で定義した LAR/ASR である．

\begin{figure}[H]
  \centering
  \includegraphics[width=\linewidth]{figures/p_reorder_dual.pdf}
  \caption{パケット並び替え率と正規受理率（LAR）および攻撃成功率（ASR）の関係
  （$p_{\mathrm{loss}} = 0.10$）。
  実線は LAR（左軸），破線は ASR（右軸，対数スケール）を示す。}
  \label{fig:p_reorder}
\end{figure}

まず，並び替えが存在しない $p_{\mathrm{reorder}} = 0$ のときには，
可用性はいずれの防御方式でもおおむね $90\%$ 前後であり，
実験1の $p_{\mathrm{loss}} = 0.10$ の結果と整合している．
この条件では，ASR は none モードで 89.6\%，rolling 0.1\%，
window 0.5\%，challenge 0.1\% となっており，
防御なしの場合のみリプレイ攻撃が高確率で成功する．

一方，$p_{\mathrm{reorder}}$ を 0.30 まで増加させると，
LAR の挙動には方式間で大きな差が現れる．
とくに rolling では LAR が 90.3\% から 76.8\% まで
13.5 ポイント低下し，最終的には正規コマンドの 23.2\%
（およそ 4 回に 1 回）が誤って拒否される状況となる．
これは，乱序が存在しない場合と比べて，
約 13.5 ポイント分（およそ 7 回に 1 回に相当）の
追加の誤拒否が発生していることを意味する．
受信側がシーケンス番号の単調増加を厳密に要求するため，
単に到着順序が前後しただけのフレームを「再生済み」と誤判定してしまうことによる．
同様に，FIDO 認証器では単調増加カウンタを用いてクローン検知を行う設計が採用されており\cite{fido-security-ref}，
乱序やロールバックが発生した場合にこれらを適切に扱えなければ，可用性が低下し得る．

これに対して window 方式では，
$p_{\mathrm{reorder}} = 0.30$ においても LAR は 90.6\% とほぼ変化せず，
ASR も 0.5\% 程度に抑えられている．
これは IPsec ESP のアンチリプレイ機能と同様に，
スライディングウィンドウとビットマップにより
「許容される乱序幅」を持つ受信窓を維持することで，
順序が前後したフレームを過剰に拒否しない設計になっているためである\cite{rfc4303,rfc6479}．

challenge 方式では，ASR は全域で 0.1\% 程度に抑えられている一方，
LAR は $p_{\mathrm{reorder}} = 0.30$ で 64.5\% まで低下し，
rolling よりも可用性への影響が大きい．
チャレンジ--レスポンスは本質的に双方向の往復を前提としており，
チャネル損失と並び替えが重なるとチャレンジと応答のペアを
正しく照合できないケースが増えるためである．
すなわち，本実験条件では challenge は
「高いセキュリティを保ったまま可用性を犠牲にする」振る舞いを示している．

以上をまとめると，実験2から次の知見が得られる．
\begin{itemize}
  \item none モードでは，乱序の有無にかかわらず ASR が約 90\% と高く，
        防御機能が無ければリプレイ攻撃は容易に成功する．
  \item rolling は乱序に対して非常に脆弱であり，
        RFC~4737 で定義されるパケット並び替えメトリクス\cite{rfc4737}で
        問題となる程度の乱序を許容する用途には不適である．
  \item window は本実験範囲の乱序に対して高い可用性と低い ASR を同時に維持しており，
        インターネット VPN などで広く採用されている
        IPsec のスライディングウィンドウ型アンチリプレイ機能と
        類似した，バランスの良い防御方式と言える\cite{rfc4303,rfc6479}．
  \item challenge は ASR をほぼゼロに抑えられるが，
        乱序と損失が重なる環境では可用性の低下が大きく，
        低遅延かつ乱序の少ないチャネル，
        あるいは高セキュリティ用途に限定して用いるのが現実的である．
\end{itemize}

\subsection{実験3：スライディングウィンドウサイズのトレードオフ分析}

\label{sec:exp3_window}

本節では，スライディングウィンドウ方式において
受理ウィンドウ幅 $W$ が可用性（正規受理率）と
セキュリティ（攻撃成功率）の両方に与える影響を評価する．
IPsec ESP の標準仕様 \cite{rfc4303} および
そのアンチリプレイ拡張 \cite{rfc6479} では，
受信側がシーケンス番号に対して幅 $W$ のスライディングウィンドウと
ビットマップを保持し，
ウィンドウ左側（古すぎるパケット）および重複パケットを拒否する方式が採用されている．
RFC~4303 は，32 ビットシーケンス番号を想定した場合に
\emph{少なくとも 32 パケット} のウィンドウサイズを実装し，
\emph{64 パケット} をデフォルト値として用いることを推奨しているが，
それ以上のウィンドウサイズについては，
想定される並び替えの程度や実装コストに応じて受信側が選択してよいと規定している
\cite{rfc4303}．

\subsubsection{実験条件}

本実験では，防御メカニズムをスライディングウィンドウ方式に固定し，
ウィンドウ幅 $W$ のみを変化させた．
他の実験と同様に，モンテカルロシミュレーションにより
確率的チャネル環境下での性能を評価する．
主な条件を以下にまとめる．

\begin{itemize}
  \item 可変パラメータ：
    \begin{itemize}
      \item 受理ウィンドウ幅 $W \in \{1, 3, 5, 7, 9, 15, 20\}$
    \end{itemize}
  \item チャネル条件（ストレスシナリオ）：
    \begin{itemize}
      \item パケット損失率 $p_{\mathrm{loss}} = 0.15$
      \item パケット並び替え率 $p_{\mathrm{reorder}} = 0.15$
      \item 損失と並び替えはいずれも独立同分布 (i.i.d.) モデルとし，
            バースト損失は未考慮とする
    \end{itemize}
  \item 攻撃モデル：
    \begin{itemize}
      \item inline 型リプレイ攻撃者（正規フレームの直後に再生フレームを挿入）
      \item 攻撃者は正規フレームを完全に記録できると仮定し，記録損失は 0 とする
    \end{itemize}
  \item モンテカルロ条件：
    \begin{itemize}
      \item 各構成について 200 回のモンテカルロ試行
      \item 1 試行あたりの正規送信数 $N_{\mathrm{legit}} = 20$
      \item 1 試行あたりのリプレイ攻撃試行数 $N_{\mathrm{attack}} = 100$
      \item 乱数シードは全構成で共有し，
            チャネル系列とコマンド列を揃えることで比較を公平化
    \end{itemize}
  \item 評価指標：
    \begin{itemize}
      \item 正規受理率（Legitimate Acceptance Rate; LAR）
      \item 攻撃成功率（Attack Success Rate; ASR）
      \item 総合スコア $J = \mathrm{LAR} - \mathrm{ASR}$
            （単純な可用性--セキュリティ折衷指標；ここで $J$ は正規受理率 LAR と攻撃成功率 ASR の差に基づく総合スコアを表す）
    \end{itemize}
\end{itemize}

実験データは，シミュレータの出力ファイル
\texttt{results/window\_sweep.json} から取得した．
図\ref{fig:window_tradeoff} に，$W$ に対する LAR・ASR・総合スコアの
関係を示す．

\begin{figure}[H]
  \centering
  \includegraphics[width=\linewidth]{figures/window_tradeoff.pdf}
  \caption{スライディングウィンドウ幅 $W$ に対する
  正規受理率・攻撃成功率・総合スコア
  （$p_{\mathrm{loss}}=0.15,\,p_{\mathrm{reorder}}=0.15$, inline 攻撃）}
  \label{fig:window_tradeoff}
\end{figure}

\subsubsection{結果と観察}

特に，$W=1$，$W=3$--$7$，$W \ge 9$ の 3 つの領域で
振る舞いが大きく異なる．

\begin{itemize}
  \item $W=1$（極端に小さいウィンドウ）：
    \begin{itemize}
      \item 正規受理率は 25.9\% まで急落し，本研究の想定ユースケースでは
            ほぼ実用的でない．
      \item 攻撃成功率は 7.3\% と中程度にとどまるが，
            可用性の大幅な低下に比べるとメリットは小さい．
      \item 総合スコア $J$ は 18.6 と，他の設定に比べて極端に低い．
    \end{itemize}
  \item $W=3$--$7$（バランスが良い領域）：
    \begin{itemize}
      \item 正規受理率はいずれも 85\% 前後で高く維持される．
      \item 攻撃成功率は $W=3$ で 6.5\%，$W=7$ でも 8.7\% にとどまり，
            総合スコア $J$ は 76.7--78.6 の高い値を示す．
      \item 本実験条件下では $W=3$ が $J=78.6$ で最良だが，
            $W=5$（$J=77.7$）および $W=7$（$J=76.7$）との差は数ポイント以内であり，
            統計的には同程度の性能と見なせる．
    \end{itemize}
  \item $W \ge 9$（過度に大きなウィンドウ）：
    \begin{itemize}
      \item 正規受理率は 85.5\% でほぼ飽和しており，
            $W$ の増加による改善は見られない．
      \item 一方で攻撃成功率は，$W=9$ で 9.6\%，$W=20$ では 11.6\% まで単調に増加し，
            総合スコア $J$ は 75.9 から 73.9 へと低下する．
      \item したがって，$W$ を 9 以上に拡大しても可用性の向上は得られず，
            セキュリティ面だけが悪化する結果となる．
    \end{itemize}
\end{itemize}

\subsubsection{考察：ウィンドウサイズ設計の指針}

以上の結果から，本シミュレーション条件
（$p_{\mathrm{loss}}=0.15,\,p_{\mathrm{reorder}}=0.15$, inline 攻撃）においては，
スライディングウィンドウ幅 $W$ の「実用的な範囲」は
概ね $W=3$～$7$ にあると結論づけられる．
この範囲では，可用性 85\% 前後と比較的低い攻撃成功率
（6.5--8.7\%）のバランスが得られる一方，
$W=1$ では可用性が崩壊し，$W \ge 9$ では攻撃成功率が有意に増加する．

IPsec ESP の Anti-Replay 仕様 \cite{rfc4303,rfc6479} では，
32～64 パケット程度のウィンドウを最低限とした上で，
\emph{より大きなウィンドウサイズは，システム要件や許容したいパケット並び替えの程度に応じて
実装者が選択する} ことが認められている \cite{rfc4303,rfc6479}．
本研究の対象である 2.4\,GHz 帯無線リモコンシステムでは，
1 セッションあたりのフレーム数が非常に小さい一方向制御であり，
IPsec VPN のような高スループット双方向トンネルとは前提条件が大きく異なる．
そのため，本研究では小さな整数値の $W$ を掃引し，
モンテカルロシミュレーションにより
\begin{itemize}
  \item 正規パケットの並び替えをどこまで許容すべきか（可用性の観点）
  \item ウィンドウの拡大に伴う攻撃面積の増加をどこまで許容できるか（セキュリティの観点）
\end{itemize}
を定量的に評価した．

その結果，本チャネル条件では $W=3$～$7$ が
「安全かつ実用的な領域」であり，
特に $W=5$ は
\begin{itemize}
  \item $W=3$ に比べて総合スコアの低下は 0.9 ポイントにとどまり，
  \item より大きな並び替え・損失が発生した場合にも
        追加のマージンを与える中庸な値である
\end{itemize}
ことから，本研究のデフォルト設定として採用した．

ターゲット環境のネットワーク特性がよく分かっており，
並び替えや損失が比較的軽微である場合には $W=3$ を，
条件が不確実または変動が大きい場合には $W=5$ を選択することで，
IPsec の標準仕様が示す設計方針 \cite{rfc4303,rfc6479} と整合的なかたちで
ウィンドウサイズをチューニングできる．

\subsection{防御メカニズムの総合比較}

\label{subsec:overall-comparison}

本節では，実験1～3の結果と各方式のアルゴリズム構造に基づき，
4種類の防御メカニズムを
\textbf{セキュリティ（攻撃成功率 $ASR$）}，
\textbf{可用性（正規受理率 $LAR$）}，
\textbf{実装コスト・要求インフラ}
の 3 観点から整理する．
数値はいずれも本章で用いたモンテカルロシミュレータ\cite{ReplaySimReadmeJP}
に基づく代表値であり，
独立損失チャネルおよび inline 型攻撃者という前提の下で得られたものである．

\vspace{1em}
\noindent\textbf{1) 防御なし方式（none）}\par
none 方式はフレームをそのまま受理するため，
例えば実験2の基準条件である $p_{\mathrm{loss}}=0.1, p_{\mathrm{reorder}}=0$ の場合，
$LAR \approx 90\%$, $ASR \approx 90\%$ と，
高い可用性と同時に極めて高い攻撃成功率を示す．
攻撃者が 1 回記録したフレームを繰り返し再生できるため，
リプレイ攻撃に対して本質的な防御機能を持たない．
本方式は実装コストが最も低く，評価用ベースラインとしては有用だが，
実システムで採用すべき選択肢ではない．

\vspace{1em}
\noindent\textbf{2) ローリングカウンタ方式（rolling）}\par
rolling 方式は単調増加するカウンタのみを受理する設計であり，
損失チャネル下では本研究の範囲で
$ASR \approx 0.1\%$ と非常に低い攻撃成功率を維持できる．
一方で，実験2が示すように $p_{\mathrm{reorder}}=0.3$ の乱序環境では，
$LAR$ が $90.3\%$ から $76.8\%$ に 13.5 ポイント低下し，
遅延して到着した正規フレームをリプレイとして誤検出してしまう．
すなわち rolling は「単純だが乱序に極端に弱い」メカニズムであり，
マルチパス伝送やリンク層再送が存在する 2.4\,GHz 帯システムとは相性が良くない．

\vspace{1em}
\noindent\textbf{3) スライディングウィンドウ方式（window）}\par
window 方式は，IPsec などで用いられるスライディングウィンドウ型アンチ
リプレイ\cite{rfc6479}を簡略化したものであり，
ある程度の並び替えを許容しつつ再生フレームのみを拒否できる．
実験2では $p_{\mathrm{reorder}}=0.3$ でも
$LAR$ は $90.3\%\rightarrow 90.6\%$ とほぼ維持され，
$ASR \approx 0.5\%$ に抑えられた．
さらに実験3では，ウィンドウ幅 $W=3$～$7$ の範囲で
$LAR \approx 85\%$, $ASR \approx 6.5$～$8.7\%$，
総合スコア $J = \mathrm{LAR} - \mathrm{ASR}$ が最大となることが確認された．
$W$ が小さすぎると可用性が（$W=1$ で $LAR=25.9\%$），
大きすぎるとセキュリティが（$W \ge 9$ で $ASR \ge 9.6\%$）悪化するため，
本研究の条件では $W=3$～$5$ 程度が実用的な折衷案と言える．

\vspace{1em}
\noindent\textbf{4) チャレンジレスポンス方式（challenge）}\par
challenge 方式は毎回新しい nonce に対する応答のみ受理するため，
実験1・2の全ての条件で攻撃成功率は常に $ASR < 0.3\%$ に抑えられ，
本研究で評価した方式の中で最も強力な防御性能を示した．
一方，チャレンジメッセージと応答の 2 段階が必要なため，
パケット損失や乱序の影響を受けやすく，
$p_{\mathrm{reorder}}=0.3$ では $LAR$ が
$89.8\%$ から $64.5\%$ まで低下した．
また，送受信双方に乱数生成器と MAC 計算機能を追加する必要があり，
片方向の単純なリモコンシステムにそのまま適用するのは負荷が大きい．

\vspace{1em}
\noindent\textbf{5) 総合評価}\par
以上をまとめると，本研究のチャネルモデルと攻撃モデルの範囲では，
次の順位付けが得られる．
\begin{itemize}
  \item \textbf{セキュリティ（ASR 単独）}：
        challenge $\gtrsim$ rolling $\approx$ window $\gg$ none
  \item \textbf{可用性とセキュリティの両立}：
        window（$W=3$～$5$）$>$ rolling $\approx$ challenge $\gg$ none
\end{itemize}

特に，乱序が無視できない無線環境では，
適切に設計されたスライディングウィンドウ方式が
単純なローリングカウンタ方式よりも実用的であり，
チャレンジレスポンス方式は
「セキュリティを最優先し，可用性をある程度犠牲にできる場面」で
有力な選択肢となる，という点が本章の重要な結論である．

\subsection{実機リモコンカーへの適用可能性の検討}

\label{subsec:rc-car-applicability}

本節では，第3章で解析した市販 2.4\,GHz 帯リモコンカーに対して，
第4章で定義した 4 種類の防御メカニズム
（防御なし（none），ローリングカウンタ方式（rolling），
スライディングウィンドウ方式（window），
チャレンジレスポンス方式（challenge））を
そのまま適用できるかどうかを検討する．
具体的には，実機プロトコルのフレーム構造とマイコン資源，
および 2.4\,GHz 帯チャネルの性質を整理したうえで，
第\ref{sec:exp-overview}節～第\ref{sec:exp3_window}節の
シミュレーション結果がどの程度まで実機に反映可能かを議論する．

\subsubsection{対象リモコンカーとチャネル条件の整理}

第3章で示したように，対象リモコンカーは
送信側リモコンから受信側車体への一方向リンクで構成されており，
ACK や暗号化されたハンドシェイクは存在しない．
フレーム構造は，先頭のアドレス＋固定ヘッダ部と，
左右輪の駆動方向や速度，灯光の ON/OFF などを表す
5 バイトの制御フィールド，末尾の CRC から構成される．
複数コマンドを比較した結果，制御フィールド末尾には
「ほぼ常に異なる」ビット列が付与されており，
第4章で述べたように，簡易的なローリングカウンタあるいは
シーケンス番号として解釈できることが分かっている．
一方で，受信側は CRC の成否以外の状態を保持しておらず，
過去に受信したことのあるフレームかどうかを判定する
メモリを一切持たないため，
記録したフレームの単純な再送であっても，
チャネル誤りがなければ正規コマンドとして受理される．
これは，433\,MHz 帯の簡易リモコンやゲート制御装置が
リプレイ攻撃に脆弱であることを示した研究
\cite{Mykhaylova2024RCReplay,FSK_RemoteAN}や，
多数の家庭用 IoT デバイスがローカル通信のリプレイ攻撃に対して
十分な保護を備えていないことを示した REPLIOT の評価
\cite{Lazzaro2024,LazzaroREPLIOT}とも整合的である．

本研究のシミュレーションでは，
README に明示したとおり\cite{ReplaySimReadmeJP}，
チャネルを損失確率 $p_{\mathrm{loss}}$ と
並び替え確率 $p_{\mathrm{reorder}}$ を持つ i.i.d.\ モデルとして扱い，
\[
  0 \leq p_{\mathrm{loss}} \leq 0.30,\qquad
  0 \leq p_{\mathrm{reorder}} \leq 0.30
\]
の範囲で掃引した．
RFC~8868 はインタラクティブリアルタイムメディアの評価指針として
0～20\% 程度の損失条件を典型例として挙げており\cite{RFC8868}，
パケット損失・乱序を対象としたモデル化や性能評価研究も
数十\% 程度までの悪条件を取り扱っている\cite{rfc4737,RFC5170,Neumann2005FEC}．
したがって，本論文で採用したパラメータ範囲は，
市販リモコンカーが混雑した 2.4\,GHz 帯チャネル上で動作する状況を含む
「現実的～やや悲観的な」条件をカバーしていると言える．
ただし，本モデルはバースト損失や時間相関を明示的には扱っておらず，
第7章で述べるように，Gilbert--Elliott 型モデル\cite{Hasslinger2008GE}などへの拡張が
今後の課題として残されている．

\subsubsection{各防御メカニズムの実装可能性}

\vspace{1em}
\noindent\textbf{1) 防御なし（none）}\par
none 方式は，現在のリモコンカーと同様に，
受信側が CRC のみを検査し，過去の受信履歴を全く保持しない構成である．
第\ref{sec:exp-loss}節および第\ref{sec:exp-reorder}節の結果
（図\ref{fig:p_loss}, 図\ref{fig:p_reorder}）から，
none 方式は理想チャネルでは LAR がほぼ 100\% であり，
可用性の観点では最も有利である一方，
攻撃者が一度でも正規通信を取得できれば，
ASR もほぼ 100\% に達してしまうことが分かる．
Anthi らによる IoT デバイスの RF 脆弱性調査\cite{Anthi2024RFVuln}や，
Lazzaro らの REPLIOT による評価\cite{Lazzaro2024,LazzaroREPLIOT}でも，
暗号学的保護やリプレイ防御を持たないデバイスは，
単純な重放だけで高確率で制御可能であることが示されており，
本研究で得られた none 方式の結果は，こうした実機観測とも一致している．

\vspace{1em}
\noindent\textbf{2) ローリングカウンタ方式（rolling）}\par
rolling 方式は，第4章で定義したように，
送信側・受信側が共有するローリングカウンタ $ctr$ と
メッセージ認証コード（MAC）を用いて，
古いフレームの再送を検出・拒否する方式である．
センサネットワーク向けの SPINS/SNEP\cite{PerrigSPINS} や，
自動車のリモートキーレスエントリ（RKE）システム\cite{Garcia2016RKE,Francillon2011PKES}では，
数バイト程度のカウンタと対称鍵ベースの MAC を組み合わせた方式が
低価格マイコン上で広く実装されている．
対象リモコンカーに rolling を導入する場合，
第3章で観測された「ほぼ常に異なる」末尾ビット列を
明示的なカウンタとして再定義し，
制御フィールドの一部を短い MAC に割り当てる，
あるいはフレーム長を数バイト拡張して
カウンタ + MAC 用の領域を新たに追加する，といった設計が考えられる．
いずれの場合も，追加で保持すべき状態は
カウンタ値と共有鍵，MAC 計算用の一時領域のみであり，
第2章および第4章で想定した低コストマイコンの RAM 制約と整合的である．
一方で，第\ref{sec:exp-reorder}節で示したように，
$p_{\mathrm{reorder}}$ が大きい環境では rolling 方式の LAR が顕著に低下し，
遅延して到着した正規フレームをリプレイとして誤って棄却してしまう点が課題となる．

\vspace{1em}
\noindent\textbf{3) スライディングウィンドウ方式（window）}\par
window 方式は，IPsec ESP におけるアンチリプレイ機能\cite{rfc4303,rfc6479}を簡略化したものであり，
受信側が「最新カウンタ値から幅 $W$ の範囲」をウィンドウとして保持し，
その範囲内で未使用のシーケンス番号のみを受理する．
実装上は rolling と同様にシーケンス番号と MAC をフレームに含める必要があるが，
ウィンドウ幅 $W$ に対応するビットマップは，
$W=3$～$7$ 程度であれば 1\,byte 程度で表現できる．
第\ref{sec:exp-loss}節，第\ref{sec:exp-reorder}節および第\ref{sec:exp3_window}節の結果から，
$W=3$～$7$ 程度の window 方式は，本研究で想定した
$p_{\mathrm{loss}} \leq 0.30,\, p_{\mathrm{reorder}} \leq 0.30$ の範囲において，
チャネルが厳しい場合でも LAR を 70～100\% 程度に維持しつつ，
post 型リプレイに対する ASR を 1\% 未満，
inline 型リプレイに対しても数\% 程度に抑えられることが分かった．
このことから，本研究で対象としたクラスの市販リモコンカーに対しては，
フレームフォーマットとファームウェアのみの変更で導入可能な範囲において，
window 方式が最も現実的な防御候補であると結論づけられる．

\vspace{1em}
\noindent\textbf{4) チャレンジレスポンス方式（challenge）}\par
challenge 方式は，受信側が乱数ノンス $N$ を送信し，
送信側が $N$ とコマンド $C$，共有鍵 $K$ に基づいて MAC を計算して応答することで，
各コマンドに対して一度きり有効なフレームのみを受理する方式である\cite{PerrigSPINS}．
第\ref{sec:exp-loss}節および第\ref{sec:exp-reorder}節の結果から，
本方式は全てのチャネル条件で ASR をほぼ 0\% まで抑えられ，
最も強力な防御性能を示す．
同様のチャレンジレスポンスに基づくリプレイ防御は，
高機能 IoT デバイス向けの認証方式としても提案されている\cite{Nain2017,Feng2017}．
しかし，チャレンジを送信する上りリンクと，
チャレンジ／応答の 2 段階通信に伴うレイテンシを許容する設計が必要であり，
現状の一方向リモコンカーに対しては，
ハードウェア構成そのものを見直さない限り導入は困難である．

\subsubsection{実機適用に向けたまとめと課題}

以上の議論を総合すると，本研究で対象としたクラスの市販 2.4\,GHz リモコンカーに対しては，
\begin{itemize}
  \item フレームフォーマットを拡張してシーケンス番号と短い MAC タグを導入し，
        受信側で幅 $W$ のスライディングウィンドウを維持する window 方式が，
        追加ハードウェアを必要とせずに実装可能な範囲で最も有望な防御方式である．
  \item rolling 方式は実装コストの観点では window と同程度であるが，
        乱序に対する耐性が低く，混雑した 2.4\,GHz 帯環境では
        可用性の面で不利になりうる．
  \item challenge 方式はセキュリティの観点では最も魅力的であるものの，
        双方向通信路とより高機能なハードウェアを前提とするため，
        本研究で扱った低コスト玩具カーよりも，
        将来の高機能 IoT デバイス向けの選択肢として位置付けられる．
\end{itemize}

なお，本研究で用いた市販トイカー自体は，
送信機・受信機ともにクローズドな RF モジュールとして実装されており，
ファームウェアやプロトコルを外部から書き換える手段は公開されていない．
したがって，本論文で提案した window 方式や challenge 方式を
当該製品に直接適用し，その場でリプレイ攻撃を無効化することは
現時点では困難である．
本節の議論は，このトイカーを典型例としつつ，
同様の構成を持つ 2.4\,GHz 帯無線リモコンシステムを新たに設計する場合に，
どの防御方式を採用すべきかを検討するための設計指針を与えることを
目的とするものである．

もっとも，これらの結論は，
i.i.d.\ 損失チャネルと inline 型攻撃者という理想化モデルに基づくものであり，
実際の 2.4\,GHz 帯電波環境やユーザ操作パターンを完全に反映したものではない．
実機への適用性をより厳密に評価するためには，
市販プロポやトランシーバモジュール\cite{RC2p4GSlides,TamiyaRCGuide,Nordic_nRF24L01}を用いた
パケットロス・遅延統計の測定や，
window 方式を実際のリモコンカーに試験実装したうえでの長期運用評価が必要である．
これらは第7章で述べる今後の課題の一部であり，
本節の検討は，実機設計に先立って「どの程度の状態管理を追加すれば，
チャネル品質を大きく損なわずにリプレイ攻撃耐性を向上できるか」を見積もるための
指針を与えるものと位置づけられる．

\subsection{本章のまとめ}

\label{subsec:exp-summary}

本章では，第4章で定義したシステムモデルと第5章で構築した
モンテカルロシミュレータ\cite{ReplaySimReadmeJP}を用いて，
防御なし（none），ローリングカウンタ方式（rolling），
スライディングウィンドウ方式（window），
チャレンジレスポンス方式（challenge）の
4 種類のリプレイ攻撃防御メカニズムを比較評価した．
第\ref{sec:exp-overview}節で述べたように，各チャネル条件
$\theta = (p_{\mathrm{loss}}, p_{\mathrm{reorder}}, W, \ldots)$ ごとに
$N_{\mathrm{runs}} = 200$ 回の独立試行を行い，
正規受理率（Legitimate Acceptance Rate; LAR）と
攻撃成功率（Attack Success Rate; ASR），
および $J = \mathrm{LAR} - \mathrm{ASR}$ という総合スコアを評価指標とした．
これは，多数の試行に基づいて誤受理率・誤拒否率を統計的に評価する
FIDO Biometrics Requirements\cite{FIDO-Bio} や，
複数の損失条件で輻輳制御アルゴリズムを比較する
RFC~8868\cite{RFC8868} に代表される，
確率量の統計的評価という枠組みと整合的である．

実験1（第\ref{sec:exp-loss}節）では，
パケット並び替えを考慮しない $p_{\mathrm{reorder}} = 0$ の条件で
損失率 $p_{\mathrm{loss}}$ を掃引し，
パケット損失が各方式に与える影響を評価した．
その結果，いずれの方式でも $p_{\mathrm{loss}}$ の増加に伴い LAR は
ほぼ単調に低下し，
$p_{\mathrm{loss}} = 0.30$ ではおおよそ 70\% 前後に収束する一方
（例：none 70.3\%，rolling 70.3\%，window 69.5\%，challenge 70.0\%），
rolling，window，challenge の 3 方式では
$p_{\mathrm{loss}} = 0.30$ においても
ASR をそれぞれ 0.4\%，1.8\%，0.3\% に抑えられることが分かった．
すなわち，本研究で想定した独立損失チャネルの範囲では，
可用性低下の主因はチャネル側のパケット損失であり，
適切に設計された防御ロジックによる追加損失は限定的であると解釈できる．

実験2（第\ref{sec:exp-reorder}節）では，
$p_{\mathrm{loss}} = 0.10$ を固定した上で
並び替え率 $p_{\mathrm{reorder}}$ を変化させ，
乱序が各方式に与える影響を評価した．
none モードでは，乱序の有無にかかわらず
ASR が約 90\% と高止まりしており，
防御機能が存在しない場合には
リプレイ攻撃が容易に成功することが定量的に示された．
rolling 方式は乱序に対して脆弱であり，
$p_{\mathrm{reorder}} = 0.30$ では LAR が 90.3\% から 76.8\% まで低下した．
一方 window 方式では，
$p_{\mathrm{reorder}} = 0.30$ においても
LAR が 90.6\% 程度でほぼ一定かつ
ASR も 0.5\% 程度に抑えられており，
IPsec ESP のスライディングウィンドウ型アンチリプレイ機能
\cite{rfc4303,rfc6479} と同様に，
一定幅の受信窓を持つ設計が
乱序とリプレイを同時に扱ううえで有効であることが確認された．
challenge 方式は，全条件で ASR をほぼ 0.1\% に抑えられる一方，
$p_{\mathrm{reorder}} = 0.30$ で LAR が 64.5\% まで低下しており，
高いセキュリティを維持する代償として可用性が犠牲になることが分かった．

実験3（第\ref{sec:exp3_window}節）では，
window 方式におけるウィンドウ幅 $W$ を変化させ，
同一チャネル条件（$p_{\mathrm{loss}} = p_{\mathrm{reorder}} = 0.15$，inline 攻撃）における
LAR，ASR，総合スコア $J$ の変化を詳細に評価した．
その結果，特に $W=1$，$W=3$--$7$，$W \ge 9$ の
3 つの領域で挙動が大きく異なることが分かった．
$W=1$ では LAR が 25.9\% まで急落し，本研究の想定ユースケースでは実用的でない．
$W=3$--$7$ の範囲では LAR が 85\% 前後で高く維持されつつ，
ASR も 6.5\%～8.7\% 程度に抑えられ，
総合スコア $J$ が最大となることが確認された．
一方 $W \ge 9$ の領域では，LAR は 85.5\% 程度で飽和しているのに対し，
ASR が 9.6\% 以上に増加しており，
ウィンドウを過度に広げるとセキュリティのみが悪化することが示された．
したがって，本研究で想定した 2.4\,GHz 帯リモコンシステムでは，
$W=3$～$5$ 程度が可用性とセキュリティのバランスが取れた実用的な設計点である．

以上の実験結果および第\ref{subsec:overall-comparison}節の整理から，
本研究で扱ったチャネルモデルと攻撃モデルの範囲では
次のような結論が得られる．
\begin{itemize}
  \item none モードは，どのチャネル条件でも ASR が高く，
        リプレイ攻撃に対してはほぼ無防備である．
  \item rolling 方式は，損失のみの環境では良好なバランスを達成するものの，
        現実的な乱序チャネルに対しては LAR が大きく低下するため，
        乱序が発生しやすいインターネット VPN やマルチホップな IoT ネットワークに
        そのまま適用するのは危険である．
  \item window 方式は，適切な $W$ を選択することで，
        損失と乱序が混在するチャネルに対しても
        高い LAR と低い ASR を両立できることが示され，
        IPsec のスライディングウィンドウ型アンチリプレイ機能\cite{rfc4303,rfc6479}
        と同様に，バランスの良い防御方式である．
  \item challenge 方式は，ASR をほぼゼロに抑えられるという意味で最も強力であるが，
        双方向のチャレンジ--レスポンスと MAC 計算を前提とするため，
        レイテンシと実装コストの観点から適用場面が限定される．
\end{itemize}

以上の結果および第\ref{sec:defense-overview}章で整理したシステム要件を踏まえると，
本研究で対象とする低コスト 2.4\,GHz 無線リモコンシステムに対しては，
次のような防御方式の選択が妥当である．

\begin{enumerate}
  \item 市販トイカーに代表される一方向・低コストの 2.4\,GHz リモコンシステムについては，
        制御フィールドにシーケンス番号と短い MAC タグを導入し，
        受信側で $W=3$～$5$ 程度のスライディングウィンドウを維持する
        window+MAC 構成を基本方式として採用する．
  \item 応答リンクが利用でき，かつドアロックやゲート制御のように
        より高いセキュリティが要求される用途については，
        チャレンジレスポンス方式を優先的に採用し，
        必要に応じてウィンドウ方式を補助的な防御として併用する．
\end{enumerate}

以上より，本研究の結論としては，
「低コスト 2.4\,GHz 無線リモコンシステムに対する実用的な第一候補は
$W=3$～$5$ のスライディングウィンドウ方式（+MAC）であり，
チャレンジレスポンス方式は双方向通信と追加リソースを許容できる
高セキュリティ用途における上位互換的な選択肢である」と位置付けられる．

これらの知見は，第2章で整理した既存研究，
すなわち RF レベルのリプレイ攻撃とジャミングが
市販 IoT デバイスに対して広く成立しうることを示した研究\cite{Anthi2024RFVuln}，
433\,MHz 帯リモコンプロトコルのリプレイ耐性を詳細に分析した研究\cite{Mykhaylova2024RCReplay}，
多数の市販 IoT 機器に対してローカルネットワーク上のリプレイ攻撃脆弱性を
自動検査する REPLIOT ツール\cite{Lazzaro2024,LazzaroREPLIOT} などとも整合的である．
これらの研究が実機デバイスにおけるリプレイ攻撃脆弱性の広がりを示しているのに対し，
本章の結果は，
固定コード型 2.4\,GHz リモコンシステムに対して
どのような防御メカニズムとパラメータ設定を導入すれば，
チャネル品質を大きく損なうことなく
リプレイ攻撃耐性を高められるかという設計指針を与えるものである．

%%%%%%%%%%%%%%%%%%%%%%%%%%%%%%%%%%%%%%%%%%%%%%%%
\section{まとめと今後の課題}

\subsection{本研究の総括}

第1章では，本研究の目的として次の3点を掲げた．
\begin{enumerate}
  \item 市販 2.4\,GHz 帯無線リモコン玩具カーを題材として，一方向無線制御システムに対するリプレイ攻撃の現実性を実機実験により確認すること．
  \item 当該システムを一般化したシステム・チャネル・攻撃モデルを定式化し，防御なし方式と複数の軽量アンチリプレイ方式を共通の枠組みで比較可能にすること．
  \item チャネル損失率 $p_{\mathrm{loss}}$ と順序入れ替え率 $p_{\mathrm{reorder}}$ を考慮したモンテカルロシミュレーションにより，各方式の正規受理率（LAR）および攻撃成功率（ASR）を評価し，低コスト 2.4\,GHz 帯無線リモコンシステムの設計指針を得ること．
\end{enumerate}

本論文では，これらの目的に対して次の成果が得られた．
\begin{enumerate}
  \item まず，ソフトウェア無線機と解析ツールを用いてリモコンフレームのビット列構造を復元し，HackRF One によるリプレイ実験を通じて，正規リモコンを用いずに加減速や進行方向，灯光など車体の挙動を完全に制御可能であることを示した（第3章）．これにより，低価格一方向無線リモコンに対するリプレイ攻撃が現実的な脅威であることを実機ベースで確認した．
  \item 次に，送信側（Sender），受信側（Receiver），リプレイ攻撃者（Attacker）と，損失率 $p_{\mathrm{loss}}$ および順序入れ替え率 $p_{\mathrm{reorder}}$ をもつ確率的チャネル $C$ から成る一般化システム・脅威モデルを定式化し，防御なし方式，ローリングカウンタ方式，スライディングウィンドウ方式，簡易チャレンジレスポンス方式の4方式を同一モデル上で比較可能にした（第4章）．
  \item さらに，チャネル条件やウィンドウ幅 $W$ を掃引するモンテカルロシミュレーションを構築し（第5章），LAR と ASR，およびその差に基づく総合スコア $J$ により各方式の安全性と可用性を評価した（第6章）．その結果，本研究が想定する損失と乱序が混在する 2.4\,GHz 帯無線環境の範囲では，ウィンドウ幅 $W=3$～$5$ 程度のスライディングウィンドウ方式が，可用性を高く保ちつつ攻撃成功率を一桁台に抑える実用的な設計候補となることを示した．
\end{enumerate}

\subsection{本研究の意義と限界}

本研究の意義は，大きく次の三点に整理できる．

\vspace{0.5em}
\noindent\textbf{1) 低コスト 2.4\,GHz 無線リモコンに対する現実的なリプレイ攻撃評価基盤の構築}\par
第 3 章では，中国製市販 2.4\,GHz リモコンカーを対象に，
実際の無線フレームを HackRF One と Universal Radio Hacker を用いて解析し，
ビット列レベルでコマンド構造を復元した\cite{URH,Mykhaylova2024RCReplay,Garcia2016RKE}．
その結果として，2.4\,GHz 帯デジタルリモコンであっても，
433\,MHz 帯の簡易リモコンや車載 RKE システムと同様に，
コマンドの単純な記録・再送によって容易にリプレイ攻撃が成立し得ることを示した\cite{Anthi2024RFVuln,Lazzaro2024}．
この実機解析結果に基づき，第 4 章ではリモコン送信機，受信機，チャネル，攻撃者を含む一般化システムモデルを定式化し，
第 5 章・第 6 章では，4 種類の防御メカニズム
（none, rolling, window, challenge）を同一のモンテカルロ実験フレームワーク上で
LAR/ASR および総合スコアを用いて比較した．
このように，本研究は「極めて単純な玩具向け無線システム」と
近年報告されている IoT／RKE 系のリプレイ攻撃研究\cite{Mykhaylova2024RCReplay,AbuAlHaija2022,Quintero2023,Anthi2024RFVuln,Lazzaro2024}
との間を橋渡しする位置付けを持ち，
低コスト無線機器におけるリプレイ攻撃リスクを定量的に示した点に意義がある．

\vspace{0.5em}
\noindent\textbf{2) 防御メカニズム選択とパラメータ設計に対する定量的な指針の提示}\par
第 6 章の実験 1--3 では，パケット損失率 $p_{\mathrm{loss}}$，
並び替え率 $p_{\mathrm{reorder}}$，
ウィンドウ長 $W$ を掃引しながら 4 つの防御方式の LAR/ASR を比較した．
その結果，none モードでは常に ASR が約 90\% と高く，
防御機能が無ければリプレイ攻撃がほぼ確実に成功する一方で，
rolling/window/challenge の 3 方式はいずれも，
none モードに比べて攻撃成功率を大幅に低減し，特に rolling と challenge 方式では
実験 1--3 の全条件で ASR が 1\% 未満に抑えられることが分かった．
特に window 方式については，$W=1$ が可用性の大幅な低下を招き，
$W \ge 9$ が攻撃成功率の増大を招くのに対し，
$W=3$--7 の領域では LAR が 85\% 前後で維持されつつ
ASR は 6.5--8.7\% 程度に抑えられることが確認された．
本研究では，IPsec のアンチリプレイ機能に関する標準仕様\cite{rfc4303,rfc6479}が
通信路の特性に応じてスライディングウィンドウ長を設計すべきであると述べている点に着目し，
これを本研究で対象とした低レートなリモコンシステムにスケーリングしたうえで，
条件が比較的安定したネットワークでは $W=3$，条件が不確実または変動が大きい環境では
$W=5$ を推奨値とする設計方針を示した．
これにより，単に「どの方式が強いか」を列挙するだけでなく，
実際に IoT デバイスを設計・運用する際の具体的なパラメータ選択に役立つ知見を提供したと言える．
これらの結果に基づく方式選択およびパラメータ設計の概要を
Table~\ref{tab:design-guideline} にまとめる．

\begin{table*}[t]
  \centering
  \caption{2.4\,GHz 一方向リモコンに対するアンチリプレイ方式の設計指針のまとめ}
  \label{tab:design-guideline}
  \footnotesize
  \begin{tabularx}{\textwidth}{l l l X}
    \hline
    想定チャネル条件 & 採用する方式 & 代表的な設定値 & LAR/ASR の傾向と注意点 \\
    \hline
    損失・乱序ともに小さい &
    防御なし（none） &
    -- &
    ASR $\approx 90\%$ と高く，リプレイ攻撃がほぼ確実に成功するため実運用では非推奨． \\

    損失中程度，乱序ほとんどなし &
    rolling counter &
    -- &
    ASR はほぼ $0\%$，LAR は none とほぼ同等．ただしカウンタ同期が必須． \\

    損失中程度，軽度の乱序あり &
    sliding window &
    $W = 3$ &
    LAR はおおよそ $85\%$ を維持しつつ，
    ASR は 6.5--8.7\% 程度に抑えられる．
    小さな状態を保持する必要がある． \\

    損失中程度，乱序・条件変動が大きい &
    sliding window &
    $W = 5$ &
    LAR はわずかに低下するが ASR がさらに低下し，
    より保守的な選択となる． \\

    双方向通信や追加往復が許容される &
    challenge &
    nonce 長さ 32\,bit 程度 &
    ASR は理論上 $0\%$ まで抑えられる一方で，
    追加ラウンドと遅延が発生し，
    リアルタイム性とのトレードオフが生じる． \\
    \hline
  \end{tabularx}
\end{table*}

\vspace{0.5em}
\noindent\textbf{3) 再現可能な軽量シミュレーションフレームワークの公開}\par
本研究で用いたシミュレータは，Python によるモジュール構成とし，
送信機・受信機・チャネル・攻撃者をクラスとして分離実装したうえで，
200 回のモンテカルロ試行を標準設定とする評価スクリプトを整備した\cite{ReplaySimReadmeJP}．
さらに，攻撃者の記録モデルやパケット損失／並び替えモデル，防御メカニズムの切り替えなどを
コマンドライン引数として指定できるよう設計している．
シミュレータ本体と実験パラメータの例は GitHub 上で公開しており\cite{ReplaySimReadmeJP}，
第三者が本論文の結果を再現し，自身のデバイスやプロトコルに合わせて
防御方式やパラメータの感度分析を行える点も，本研究の実践的な意義である．

\vspace{1em}
一方で，本研究には次のような限界が存在する．

\vspace{0.5em}
\noindent\textbf{1) チャネルモデルの単純化}\par
本研究のシミュレータは，パケット損失と順序入れ替えをそれぞれ独立な確率事象として扱う
簡略化チャネルモデルに基づいている．
このモデルは i.i.d.\ 損失と遅延キューによる乱序を表現するには十分である一方で，
現実のインターネットや無線 LAN で頻繁に観測されるバースト損失や
長い尾を持つ遅延分布を再現できない\cite{Jiang2000Modeling,Hasslinger2008GE}．
実際，インタラクティブメディア向けの混雑制御評価ガイドライン\cite{RFC8868,IETF-RMCAT-Eval}や，
パケット並び替えメトリクス\cite{rfc4737}では，
Gilbert--Elliott 型の 2 状態マルコフ損失モデルや，
測定トレースに基づく複合的なチャネルモデルの使用が推奨されている．
したがって，本研究で得られた数値結果は，
「同一の簡略モデルを共有する防御方式間の相対比較」としては有効であるものの，
特定の環境における絶対的性能値としては慎重に解釈する必要がある．

\vspace{0.5em}
\noindent\textbf{2) 攻撃者モデルおよび脅威モデルの限定性}\par
本研究では，NIST CSRC の用語定義\cite{NISTReplay}に従い，
攻撃者は「合法フレームを完全に記録し，後から再送できるが，
MAC の偽造や暗号解析は行わない」という
理想化された記録・リプレイ攻撃者としてモデル化した．
しかし，実際の IoT 環境では，妨害電波によるジャミングや，
周波数ホッピングパラメータの逆解析，複数チャネルを跨いだ多点同時攻撃など，
より高度な攻撃も想定される\cite{Anthi2024RFVuln,Lazzaro2024}．
また，本研究では物理的侵入やサイドチャネル攻撃，鍵管理の欠陥といった
上位層の脅威は明示的には扱っていない．
したがって，本論文の評価は，
あくまで「無線リンク層における記録・リプレイ攻撃」に焦点を絞ったものであり，
総合的なセキュリティ評価の一部を構成する位置付けにとどまる．

\vspace{0.5em}
\noindent\textbf{3) 評価対象システムおよびパラメータ空間の限定性}\par
本研究の実機評価は特定の 2.4\,GHz リモコンカー 1 機種のみに基づいており，
またシミュレーションも，$p_{\mathrm{loss}}, p_{\mathrm{reorder}} \le 0.30$，
ウィンドウ長 $W \in \{1,3,5,7,9,15,20\}$ といった
比較的狭いパラメータ範囲に限定されている．
Wi\mbox{-}Fi，BLE，Zigbee などの標準化プロトコルや，
産業用無線，超広帯域（UWB）など他の物理層に対しても，
同様のフレームワークに基づく評価を拡張することで，
本研究の知見がどこまで一般化可能かを検証することが今後の課題である\cite{ElHajj2019IoTAuth,Anthi2024RFVuln,Lazzaro2024}．

\vspace{0.5em}
\noindent\textbf{4) 評価指標とコストモデルの簡略化}\par
本研究では，評価指標として LAR/ASR および
単純な線形結合による総合スコア $J = \mathrm{LAR} - \mathrm{ASR}$ を採用し，
計算量・メモリ使用量・消費電力といったコスト要因は明示的には評価していない．
しかし，TinySec や FlexiSec，SPINS などのセンサネットワーク向けセキュリティアーキテクチャ\cite{TinySec,JinwalaFlexiSec,PerrigSPINS}では，
鍵長や MAC 長，ウィンドウサイズの選択がノード寿命に直結することが繰り返し指摘されている．
本研究で扱った window/challenge 方式も，
組み込み実装における計算量やメモリフットプリントを考慮したうえで，
セキュリティとリソース制約のバランスを評価する必要がある．
FIDO Alliance による生体認証プロファイル\cite{FIDO-Bio,fido-security-ref}や
IoT 認証方式のサーベイ\cite{ElHajj2019IoTAuth}が示すように，
実システムの設計では「攻撃成功率だけでなくユーザビリティや運用コストを含めた多目的最適化」が要求される．

\vspace{0.5em}
\noindent\textbf{5) シミュレーション中心の評価であること}\par
本研究では，第 3 章で市販リモコンカーを用いたリプレイ攻撃の実機デモを行ったものの，
防御機能を実装したファームウェアを実機に組み込み，
長期運用下で評価するところまでは至っていない．
これは，対象とした市販トイカーの RF モジュールがクローズドであり，
外部からプロトコルやファームウェアを改修する手段が提供されていないためである．
モンテカルロシミュレーションは，広いパラメータ空間を効率的に探索できる一方で，
モデル化の妥当性に評価結果が強く依存する\cite{Raychaudhuri2008MC,Law2015Simulation,Kleijnen1995VV,Sargent2013VVA}．
今後は，本研究で得られたパラメータ指針（特に window 方式の $W$ 設計や
challenge 方式のノンス長設計）を実機プロトタイプに実装し，
フィールド試験やハードウェアインザループシミュレーションを通じて検証することが重要である．

\subsection{今後の課題}

\label{subsec:future-work}

前節で述べたように，本研究で得られた結論は，
(i) i.i.d.\ パケット損失と簡略化した乱序モデル，
(ii) 受信機付近に常駐する inline 型攻撃者，
(iii) ソフトウェアシミュレータ上での計算量評価
という前提に基づいている．
より実機に近い設計指針を得るためには，以下のような拡張が必要である．

\vspace{0.5em}
\noindent\textbf{1) 高度なチャネルモデルと実環境トレースに基づく評価}\par
既存の測定研究では，インターネットや無線ネットワークにおける
パケット損失や遅延は独立ではなくバースト的に発生することが報告されており，
2 状態マルコフ連鎖に基づく Gilbert--Elliott (GE) モデルなどが
広く用いられている\cite{Jiang2000Modeling,Hasslinger2008GE}．
また，IETF のリアルタイムメディア評価指針や IPPM メトリックでは，
トレースに基づく損失・遅延・再順序化のモデル化と評価方法が整理されている
\cite{RFC8868,IETF-RMCAT-Eval,rfc4737}．

本研究では解析の容易さを優先して i.i.d.\ 損失モデルを採用したが，
今後は市販プロポや 2.4\,GHz トランシーバモジュール
\cite{RC2p4GSlides,TamiyaRCGuide,Nordic_nRF24L01}から取得した
実環境トレースを用いてチャネルパラメータを推定し，
GE モデルなどの時間相関を持つチャネルモデルの下で
window 方式や rolling 方式の耐性を再評価する必要がある．
これにより，「ウィンドウ幅 $W$ をどの程度まで拡大しても
可用性を保てるか」といった設計判断を，
実際の 2.4\,GHz 帯環境に近い条件のもとで検証できると期待される．

\vspace{0.5em}
\noindent\textbf{2) 実機マイコン実装に基づくオーバヘッド評価}\par
本論文では，防御メカニズムの複雑さを主に理論的な計算量と
シミュレーション時間で評価したが，
実際の低コストマイコン上での実行時間・消費電力・メモリ消費量は
明示的には測定していない．
一方，センサネットワーク向けの軽量セキュリティ機能である
TinySec, SPINS, FlexiSec などの先行研究では，
8/16\,bit マイコン上に実装したうえで，
フラッシュ使用量，RAM 使用量，レイテンシ，エネルギ消費といった
詳細なオーバヘッド評価が行われている
\cite{TinySec,PerrigSPINS,JinwalaFlexiSec}．
さらに，DTLS や IPsec などの実プロトコルでは，
スライディングウィンドウ型の anti-replay 機能と暗号処理を組み合わせた
実装例が多数報告されている\cite{RFC6347DTLS,rfc4303}．

今後は，rolling 方式，window 方式，challenge 方式を
典型的な 32\,bit クラスのマイコン（例：ARM Cortex-M 系）に実装し，
上述した指標に基づくオーバヘッド評価を行う必要がある．
また，IoT 認証方式やセキュリティアーキテクチャに関する近年のサーベイ
\cite{ElHajj2019IoTAuth,Dritsas2025IoTSecurity}を踏まえ，
要求されるセキュリティ水準と利用可能な計算資源とのバランスを考慮した
実装指針を検討することが重要である．

\vspace{0.5em}
\noindent\textbf{3) 攻撃者モデルと防御メカニズムの拡張}\par
本論文では，受信機近傍に常駐し固定コード信号を完全復号できる
inline 型攻撃者という単純化したモデルを採用した．
しかし，自動車の PKES や RKE を対象とした既存研究では，
物理層リレー攻撃\cite{Francillon2011PKES}やローリングコード方式の脆弱性
\cite{Garcia2016RKE}，さらには深層学習を用いたリプレイ信号検知
\cite{AbuAlHaija2022}など，より多様な攻撃・防御が報告されている．
また，433\,MHz 帯固定コードリモコンに対する実験的研究では，
EV1527 などのプロトコルが静的コードのためリプレイ攻撃に非常に弱いことが
示されている\cite{Mykhaylova2024RCReplay}．
IoT デバイス全般に対する RF 脆弱性調査も，
ジャミングやなりすましと組み合わせた攻撃が実用的であることを指摘している
\cite{Anthi2024RFVuln}．

今後は，ジャミング＋リプレイや部分的ななりすまし攻撃を含む複合攻撃モデルを導入し，
「状態管理のみでどこまで防げるか」，「暗号学的認証や物理層特徴量と組み合わせた場合に
どのようなトレードオフが生じるか」を評価する必要がある．
これにより，本研究で扱った状態管理ベースの手法が，
どの程度まで現実の攻撃シナリオに耐えうるかを明確化できる．

\vspace{0.5em}
\noindent\textbf{4) ネットワークシミュレータとの連携と大規模評価}\par
本研究のシミュレータは単一リンク上のフレーム送受信を対象としており，
マルチホップ経路やバックホールとの相互作用はモデル化していない．
一方，ns-3 や OMNeT++ などのネットワークシミュレータは，
無線 PHY/MAC，ルーティングプロトコル，アプリケーションを含む
大規模ネットワークのエンド・ツー・エンド性能評価に広く用いられている
\cite{ns3Tutorial,OMNeTppIntro}．
IETF の評価指針\cite{RFC8868}でも，
トレース駆動型のチャネルモデルとネットワーク全体の混雑状態を
組み合わせた評価の重要性が指摘されている．

今後は，本研究のリプレイ攻撃モデルと防御アルゴリズムを
ns-3 や OMNeT++ のモジュールとして実装し，
\cite{ReplaySimReadmeJP}で定義したパラメータ空間を保ちながら，
マルチホップ IoT ネットワークにおける遅延・スループット・再送制御との
相互作用を明らかにすることが望まれる．

\vspace{0.5em}
以上のように，本研究で用いた抽象モデルとモンテカルロ評価の利点を維持しつつ，
チャネルモデル，攻撃者モデル，実機マイコン実装，ネットワーク全体の観点から段階的に現実性を高めることで，
「どのような環境・機器・攻撃シナリオに対して，どの程度の防御メカニズムとオーバヘッドが必要か」を
体系的に明らかにしていくことが今後の課題である．

%%%%%%%%%%%%%%%%%%%%%%%%%%%%%%%%%%%%%%%%%%%%%%%%
% 参考文献（正文での引用順に並べ替え）
\begin{thebibliography}{99}

\bibitem{Sinha2025IoTAnalytics}
S.~Sinha,
``State of IoT 2025: Number of connected IoT devices growing 14\% to 21.1 billion globally,''
IoT Analytics Research Blog, Oct. 2025. (online),
available at: \url{https://iot-analytics.com/number-connected-iot-devices/},
(accessed 2025-11-26).

\bibitem{Carpenter2025IEEE802154}
T.~Carpenter,
``IEEE 802.15.4: A Wireless IoT Powerhouse for Network Managers,''
CWNP Blog, Apr. 2025. (online),
available at: \url{https://www.cwnp.com/ieee-802.15.4-wireless-iot-powerhouse-network-managers/},
(accessed 2025-11-26).

\bibitem{BLE2025Sensors}
N.~Koulouras et al.,
``Evolution of Bluetooth Technology: BLE in the IoT Ecosystem,''
Sensors, vol.~25, no.~4, 996, 2025.

\bibitem{Dritsas2025IoTSecurity}
E.~Dritsas et al.,
``A Survey on Cybersecurity in IoT,''
Future Internet, vol.~17, no.~1, 30, 2025.

\bibitem{Francillon2011PKES}
A.~Francillon, B.~Danev, and S.~Capkun,
``Relay Attacks on Passive Keyless Entry and Start Systems in Modern Cars,''
Proc.\ NDSS 2011, San Diego, CA, USA, 2011.

\bibitem{Lin2019IndustrialRemote}
Y.-L.~Lin et al.,
``A Security Analysis of Radio Remote Controllers for Industrial Applications,''
Black Hat Asia, Whitepaper, Singapore, 2019.

\bibitem{Garcia2016RKE}
F.~D.~Garcia, D.~Oswald, T.~Kasper, and P.~Pavlidès,
``Lock It and Still Lose It -- On the (In)Security of Automotive Remote Keyless Entry Systems,''
in \textit{Proc. 25th USENIX Security Symposium (USENIX Security 2016)},
Austin, TX, USA, pp.~929--944, Aug. 2016.

\bibitem{Anthi2024RFVuln}
E.~Anthi et al.,
``Investigating Radio Frequency Vulnerabilities in the Internet of Things (IoT),''
IoT, vol.~5, no.~2, 18, 2024.

\bibitem{Mykhaylova2024RCReplay}
O.~Mykhaylova, A.~Stefankiv, T.~Nakonechny, T.~Fedynyshyn, and V.~Sokolov,
``Resistance to Replay Attacks of Remote Control Protocols using the 433\,MHz Radio Channel,''
in \textit{Proc. Cybersecurity Providing in Information and Telecommunication Systems (CPITS 2024)},
CEUR Workshop Proc., vol.~3654, pp.~98--110, Feb. 2024.
(online), available at: \url{https://ceur-ws.org/Vol-3654/paper9.pdf},
(accessed 2025-11-26).

\bibitem{RFC9827}
V.~Smyslov,
``Renaming the Extended Sequence Numbers (ESN) Transform Type in the Internet Key Exchange Protocol Version 2 (IKEv2),''
RFC 9827, IETF, Nov. 2025.

\bibitem{rfc3552}
E.~Rescorla and B.~Korver,
``Guidelines for Writing RFC Text on Security Considerations,''
RFC 3552, IETF, July 2003.

\bibitem{Awaad2014IPsecWindow}
M.~H.~Awaad,
``The Use of Dynamic Sliding Window with IPSec,''
\textit{Journal of Education for Pure Science},
College of Education for Pure Sciences, Univ.\ of Thi-Qar, vol.~4, pp.~278--289, 2014.
(online), available at: \url{https://iasj.rdd.edu.iq/journals/uploads/2024/12/07/b6a7fbae19e08ede0116425d8166b3c2.pdf},
(accessed 2025-11-26).

\bibitem{Serena2025SimCyber}
L.~Serena, G.~D'Angelo, S.~Ferretti, and M.~Marzolla,
``Simulation in Cybersecurity: Understanding Techniques, Applications, and Goals,''
arXiv preprint arXiv:2508.06106, Aug. 2025.
(online), available at: \url{https://arxiv.org/abs/2508.06106},
(accessed 2025-11-26).

\bibitem{RC2p4GSlides}
FTX,
``FTX Banzai RTR Brushed 2.4\,GHz 4WD Drift Car: Instruction Manual,''
CML Distribution, 2013. (online),
available at: \url{https://www.cmldistribution.co.uk/diagrams/FTX/Banzai%20Manual.pdf},
(accessed 2025-11-26).

\bibitem{TamiyaRCGuide}
タミヤ,
「タミヤRCシステム ファインスペック2.4G 2チャンネルRCプロポセット 取扱説明書」,
田宮模型株式会社, 2021. (online),
available at: \url{https://d7z22c0gz59ng.cloudfront.net/cms/japan/download/rcmanual/45053_2021.pdf},
(accessed 2025-11-26).

\bibitem{EURORC2025}
EuroRC,
``Beginner's Guide to RC Cars,''
EuroRC News, Apr. 28, 2025. (online),
available at: \url{https://www.eurorc.com/news/318/rc-car-beginner-guide},
(accessed 2025-11-26).

\bibitem{IoTAnalytics2024}
S.~Sinha,
``State of IoT 2024: Number of Connected IoT Devices Growing 13\% to 18.8 Billion Globally,''
IoT Analytics Research Blog, 2024. (online),
available at: \url{https://iot-analytics.com/number-of-connected-iot-devices-2024/},
(accessed 2025-11-26).

\bibitem{ABIWireless2024}
ABI Research,
``How Multi-Protocol Wireless Platforms Are Enabling the Next Wave of Wireless Innovation,''
Whitepaper commissioned by CEVA, 2024. (online),
available at: \url{https://www.ceva-ip.com/wp-content/uploads/ABI-Research-Protocol-Wireless-Platforms-Are-Enabling-The-Next-Wave-Of-Wireless-Innovation-Whitepaper.pdf},
(accessed 2025-11-26).

\bibitem{BLEMarket2024}
360 Research Reports,
``Bluetooth Low Energy (BLE) Modules Market Size, Share, Forecast 2034,''
Market Research Report, base year 2024. (online),
available at: \url{https://www.360researchreports.com/market-reports/bluetooth-low-energy-modules-market-200160},
(accessed 2025-11-26).

\bibitem{EN300328}
ETSI,
``EN 300 328: Wideband transmission systems; Data transmission equipment operating in the 2.4\,GHz ISM band and using wide band modulation techniques,''
European Telecommunications Standards Institute. (online),
available at: \url{https://www.etsi.org/deliver/etsi_en/300300_300399/300328/02.02.02_60/en_300328v020202p.pdf},
(accessed 2025-11-26).

\bibitem{NXP_Coexistence}
NXP Laboratories UK,
``Co-existence of IEEE 802.15.4 at 2.4\,GHz,''
Application Note JN-AN-1079 v1.1, Nov. 2013. (online),
available at: \url{https://www.nxp.com/docs/en/application-note/JN-AN-1079.pdf},
(accessed 2025-11-26).

\bibitem{Nordic_nRF24L01}
Nordic Semiconductor,
``nRF24L01+ Single Chip 2.4\,GHz Transceiver, Product Specification v1.0.'' (online),
available at: \url{https://cdn.sparkfun.com/assets/3/d/8/5/1/nRF24L01P_Product_Specification_1_0.pdf},
(accessed 2025-11-26).

\bibitem{Sklar_DC}
B.~Sklar,
\textit{Digital Communications: Fundamentals and Applications}, 2nd ed.,
Prentice Hall, 2001.

\bibitem{FSK_RemoteAN}
Linx Technologies,
``The Basics of Remote Control and Remote Keyless Entry,''
Application Note AN-00320, Rev. 11/06/2012. (online),
available at: \url{https://www.mouser.com/catalog/additional/Linx%20Technologies_an-00320.pdf},
(accessed 2025-11-26).

\bibitem{RKE_Overview}
Maxim Integrated,
``Requirements of Remote Keyless Entry (RKE) Systems,''
Application Note AN3395, Feb. 2005. (online),
available at: \url{https://www.analog.com/en/resources/technical-articles/requirements-of-remote-keyless-entry-systems.html},
(accessed 2025-11-26).

\bibitem{NISTNonce}
NIST,
``Nonce,'' NIST Computer Security Resource Center (CSRC) Glossary. (online),
available at: \url{https://csrc.nist.gov/glossary/term/nonce},
(accessed 2025-11-26).

\bibitem{SenWSNSurvey}
J.~Sen,
``A Survey on Wireless Sensor Network Security,''
International Journal of Communication Networks and Information Security,
vol.~1, no.~2, pp.~59--82, 2009.

\bibitem{Gesteira2025}
R.~Gesteira-Miñarro, G.~López, and R.~Palacios,
``Revisiting Wireless Cyberattacks on Vehicles,''
Sensors, vol.~25, no.~8, 2605, 2025.

\bibitem{rfc6479}
T.~Tsou and X.~Zhang,
``IPsec Anti-Replay Algorithm without Bit Shifting,''
RFC~6479, Internet Engineering Task Force, Jan.\ 2012.

\bibitem{Quintero2023}
J.~C.~M.~Quintero, E.~P.~E.~Cuesta, and L.~J.~R.~López,
``A new method for the detection and identification of the replay attack on cars using SDR technology and classification algorithms,''
Results in Engineering, vol.~19, 101243, 2023.

\bibitem{AbuAlHaija2022}
Q.~Abu~Al-Haija and A.~A.~Alsulami,
``Detection of Fake Replay Attack Signals on Remote Keyless Controlled Vehicles Using Pre-Trained Deep Neural Network,''
Electronics, vol.~11, no.~20, 3376, 2022.

\bibitem{Lazzaro2024}
S.~Lazzaro, V.~De~Angelis, A.~M.~Mandalari, and F.~Buccafurri,
``Is Your Kettle Smarter Than a Hacker? A Scalable Tool for Assessing Replay Attack Vulnerabilities on Consumer IoT Devices,''
in Proc.\ IEEE Int.\ Conf.\ on Pervasive Computing and Communications (PerCom), pp.~114--124, 2024.

\bibitem{LazzaroREPLIOT}
S.~Lazzaro, V.~De~Angelis, A.~M.~Mandalari, and F.~Buccafurri,
``REPLIOT: A Scalable Tool for Assessing Replay Attack Vulnerabilities on Consumer IoT Devices,''
in Proc.\ IEEE PerCom Workshops, 2024.

\bibitem{PohlNoackWOOT18}
J.~Pohl and A.~Noack,
``Universal Radio Hacker: A Suite for Analyzing and Attacking Stateful Wireless Protocols,''
in \textit{Proc.\ 12th USENIX Workshop on Offensive Technologies (WOOT~18)},
Baltimore, MD, USA, Aug.\ 2018.

\bibitem{PohlNoackIoTSP17}
J.~Pohl and A.~Noack,
``Universal Radio Hacker: A Suite for Wireless Protocol Analysis,''
in \textit{Proc.\ 2017 Workshop on Internet of Things Security and Privacy (IoTS\&P~'17)},
pp.~59--60, Dallas, TX, USA, Nov.\ 2017.

\bibitem{LinCostello2004}
S.~Lin and D.~J.~Costello,
\textit{Error Control Coding: Fundamentals and Applications},
2nd ed., Pearson (Prentice Hall), 2004.

\bibitem{KoopmanChakravarty2004}
P.~Koopman and T.~Chakravarty,
``Cyclic Redundancy Code (CRC) Polynomial Selection for Embedded Networks,''
in \textit{Proc.\ International Conference on Dependable Systems and Networks (DSN 2004)},
pp.~145--154, Florence, Italy, Jun.--Jul.\ 2004.

\bibitem{Nain2017}
A.~K.~Nain, V.~Kushwah, and R.~Kumar,
``A Secure Phase-Encrypted IEEE 802.15.4 Transceiver Design,''
IEEE Transactions on Computers, vol.~66, no.~8, pp.~1310--1323, 2017.

\bibitem{Feng2017}
Y.~Feng, W.~Wang, Y.~Weng, and H.~Zhang,
``A Replay-Attack Resistant Authentication Scheme for the Internet of Things,''
in Proc.\ 2017 IEEE Int.\ Conf.\ on Computational Science and Engineering (CSE)
and Embedded and Ubiquitous Computing (EUC), vol.~1, pp.~541--547, 2017.

\bibitem{URH}
J.~Pohl and A.~Noack,
``Universal Radio Hacker: A Suite for Wireless Protocol Analysis,''
in Proc.\ 2017 ACM Workshop on Internet of Things Security and Privacy (IoTS\&P~'17),
pp.~59--60, Dallas, TX, USA, Nov. 2017.

\bibitem{NISTReplay}
National Institute of Standards and Technology (NIST),
``Replay attack,''
NIST Computer Security Resource Center (CSRC) Glossary,
(online), available at:
\url{https://csrc.nist.gov/glossary/term/replay_attack},
(accessed 2025-11-26).

\bibitem{FIDO}
FIDO Alliance,
``FIDO Biometric Requirements, Version 3.0,''
Jan. 11, 2023.
(online), available at:
\url{https://fidoalliance.org/specs/biometric/requirements/Biometrics-Requirements-v3.0-fd-20230111.html},
(accessed 2025-11-26).

\bibitem{Jiang2000Modeling}
W.~Jiang and H.~Schulzrinne,
``Modeling of packet loss and delay and their effect on real-time multimedia service quality,''
in \textit{Proc. of NOSSDAV 2000 (Network and Operating System Support for Digital Audio and Video)},
2000.

\bibitem{Hasslinger2008GE}
G.~Ha\ss linger and O.~Hohlfeld,
``The Gilbert--Elliott model for packet loss in real time services on the Internet,''
in \textit{Proc.\ 14th GI/ITG Conf.\ on Measurement, Modelling and Evaluation of Computer and Communication Systems (MMB)},
2008, pp.~269--286.

\bibitem{rfc4737}
A.~Morton, G.~Ramachandran, S.~Shalunov, L.~Ciavattone, and J.~Perser,
``Packet Reordering Metrics,''
RFC~4737, Internet Engineering Task Force, Nov.\ 2006.

\bibitem{TinySec}
C.~Karlof, N.~Sastry, and D.~Wagner,
``TinySec: A link layer security architecture for wireless sensor networks,''
in \textit{Proc. 2nd Int. Conf. on Embedded Networked Sensor Systems (SenSys 2004)},
pp.~162--175, Baltimore, MD, USA, Nov.\ 2004.

\bibitem{JinwalaFlexiSec}
D.~C.~Jinwala, D.~R.~Patel, and K.~S.~Dasgupta,
``FlexiSec: A configurable link layer security architecture for wireless sensor networks,''
\textit{Journal of Information Assurance and Security},
vol.~4, pp.~582--603, 2009.

\bibitem{PerrigSPINS}
A.~Perrig, R.~Szewczyk, V.~Wen, D.~Culler, and J.~D.~Tygar,
``SPINS: Security Protocols for Sensor Networks,''
in \textit{Proc.\ 7th Annual Int.\ Conf.\ on Mobile Computing and Networking (MobiCom 2001)},
pp.~189--199, ACM, Rome, Italy, 2001.

\bibitem{RFC2104HMAC}
H.~Krawczyk, M.~Bellare, and R.~Canetti,
``HMAC: Keyed-Hashing for Message Authentication,''
RFC~2104, Internet Engineering Task Force, 1997.

\bibitem{RFC6347DTLS}
E.~Rescorla and N.~Modadugu,
``Datagram Transport Layer Security Version 1.2,''
RFC~6347, Internet Engineering Task Force, Jan.\ 2012.

\bibitem{rfc4303}
S.~Kent,
``IP Encapsulating Security Payload (ESP),''
RFC~4303, Internet Engineering Task Force, Dec.\ 2005.

\bibitem{ISO9798-2}
ISO/IEC 9798-2:2008,
``Information technology -- Security techniques -- Entity authentication -- Part 2: Mechanisms using symmetric encipherment algorithms'',
International Organization for Standardization, 2008.

\bibitem{mcgrew-iv-gen-03}
D.~McGrew,
``Generation of Deterministic Initialization Vectors (IVs) and Nonces'',
Internet-Draft draft-mcgrew-iv-gen-03, Work in Progress,
Internet Engineering Task Force, Oct.\ 2013.
available at: \url{https://datatracker.ietf.org/doc/draft-mcgrew-iv-gen/03/}

\bibitem{ElHajj2019IoTAuth}
M.~El-Hajj, A.~Fadlallah, M.~Chamoun, and A.~Serhrouchni,
``A Survey of Internet of Things (IoT) Authentication Schemes'',
\textit{Sensors}, vol.~19, no.~5, article 1141, 2019.

\bibitem{Raychaudhuri2008MC}
S.~Raychaudhuri,
``Introduction to Monte Carlo Simulation,''
in \textit{Proc.\ 2008 Winter Simulation Conference (WSC 2008)},
Miami, FL, USA, Dec. 2008, pp.~91--100.

\bibitem{Law2007Simulation}
A.~M.~Law,
\textit{Simulation Modeling and Analysis}, 4th~ed.,
McGraw-Hill, Boston, MA, USA, 2007.

\bibitem{Neumann2005FEC}
C.~Neumann, V.~Roca, A.~Francillon, and D.~Furodet,
``Impacts of packet scheduling and packet loss distribution on FEC performances: observations and recommendations,''
in \textit{Proc.\ 2005 ACM Conference on Emerging Network Experiment and Technology (CoNEXT~'05)},
Toulouse, France, Oct. 2005, pp.~166--176.

\bibitem{OMNeTppIntro}
OpenSim~Ltd.,
``What is OMNeT++?''
OMNeT++ official website (online).
Available at: \url{https://omnetpp.org/intro/},
accessed 2025-11-26.

\bibitem{ns3Tutorial}
ns-3~Consortium,
``ns-3 Tutorial,''
ns-3 official documentation (online).
Available at: \url{https://www.nsnam.org/docs/tutorial/singlehtml/index.html},
accessed 2025-11-26.

\bibitem{RFC5170}
C.~Neumann, V.~Roca, and D.~Furodet,
``Low Density Parity Check (LDPC) Staircase and Triangle Forward Error Correction (FEC) Schemes,''
RFC~5170, Internet Engineering Task Force, June 2008.

\bibitem{Law2015Simulation}
A.~M.~Law,
\textit{Simulation Modeling and Analysis}, 5th~ed.,
McGraw-Hill Education, 2015.

\bibitem{Sargent2013VVA}
R.~G.~Sargent,
``Verification and validation of simulation models,''
\textit{Journal of Simulation}, vol.~7, no.~1, pp.~12--24, 2013.

\bibitem{Kleijnen1995VV}
J.~P.~C.~Kleijnen,
``Verification and validation of simulation models,''
\textit{European Journal of Operational Research}, vol.~82, no.~1, pp.~145--162, 1995.

\bibitem{ReplaySimReadmeJP}
M.~Lu,
``IoT Replay Attack Defense Simulator,''
GitHub repository, 2025.
[Online]. Available:
\url{https://github.com/tammakiiroha/IoT-Replay-Defense-Simulator}

\bibitem{FIDO-Bio}
FIDO Alliance,
``FIDO Biometrics Requirements,''
Final Document, Version~4.0, May~22, 2024.
[Online]. Available:
\url{https://fidoalliance.org/specs/biometric/requirements/Biometrics-Requirements-v4.0-fd-20240522.html}

\bibitem{RFC8868}
V.~Singh, J.~Ott, and S.~Holmer,
``Evaluating Congestion Control for Interactive Real-Time Media,''
RFC~8868, Jan.~2021.
[Online]. Available: \url{https://www.rfc-editor.org/rfc/rfc8868.html}

\bibitem{fido-security-ref}
FIDO Alliance,
``FIDO Security Reference,''
FIDO Alliance Proposed Standard, Version 1.1, Feb.\ 2017.
(online), available at: \url{https://fidoalliance.org/specs/fido-uaf-v1.1-ps-20170202/fido-security-ref-v1.1-ps-20170202.html},
(accessed 2025-11-27).

\bibitem{IETF-RMCAT-Eval}
S.~Sarker, V.~Singh, X.~Zhu, and M.~Ramalho,
``Test Cases for Evaluating RMCAT Proposals,''
Internet Engineering Task Force, RFC~8867, Jan. 2021.
\url{https://www.rfc-editor.org/rfc/rfc8867}

\end{thebibliography}

\end{document}
